\documentclass[11pt,a4paper]{article}

% ─── Packages ───
\usepackage[utf8]{inputenc}
\usepackage[T1]{fontenc}
\usepackage{amsmath,amssymb,amsthm}
\usepackage{mathtools}
\usepackage{mathrsfs}
\usepackage{geometry}
\geometry{margin=1in}
\usepackage{enumitem}
\usepackage{hyperref}
\hypersetup{colorlinks=true,linkcolor=blue!60!black,citecolor=blue!60!black,urlcolor=blue!60!black}
\usepackage{booktabs}
\usepackage{fancyhdr}
\usepackage{xcolor}

% ─── Theorem environments ───
\newtheorem{theorem}{Theorem}[section]
\newtheorem{lemma}[theorem]{Lemma}
\newtheorem{corollary}[theorem]{Corollary}
\newtheorem{proposition}[theorem]{Proposition}
\newtheorem{conjecture}[theorem]{Conjecture}
\theoremstyle{definition}
\newtheorem{definition}[theorem]{Definition}
\newtheorem{assumption}[theorem]{Assumption}
\newtheorem{construction}[theorem]{Construction}
\newtheorem{axiom}{Axiom}
\theoremstyle{remark}
\newtheorem{remark}[theorem]{Remark}

% ─── Status notice ───
\newcommand{\statusnotice}[1]{%
  \medskip\noindent\textcolor{darkgray}{\textit{Status Notice.} #1}\medskip
}

% ─── Claim governance macros ───
\newcommand{\claimtag}[1]{\textcolor{blue!60!black}{\textbf{[Claim class: #1]}}}
\newcommand{\assumptionfooter}[3]{%
  \par\smallskip\noindent\textcolor{darkgray}{\footnotesize\textit{Provenance.} Depends on #1; Layer: #2; Status: #3.}\par\medskip
}

% ─── Headers ───
\pagestyle{fancy}
\fancyhf{}
\fancyhead[R]{\textit{\small Nested Fibrational Cosmology --- Canonical Monograph}}
\fancyfoot[C]{\thepage}
\renewcommand{\headrulewidth}{0.4pt}

% ─── Title ───
\title{%
  \textbf{Nested Fibrational Cosmology}\\[6pt]
  \large A Structural Derivation of Physical Invariants\\
  from Finite Relational Incidence\\[12pt]
  \normalsize Canonical Monograph --- Consolidated Edition
}
\author{Evan Blink}
\date{February 2026}

\begin{document}

\maketitle
\thispagestyle{empty}

% ═══════════════════════════════════════════
% ABSTRACT
% ═══════════════════════════════════════════
\begin{abstract}
We present a foundational framework in which physically meaningful structure is identified with invariants of finite relational incidence under a maximal class of admissible extensions. We prove that irreversible compression to survivor classes is unavoidable; structural entropy is monotone under admissible extension; and admissible laws are algebraically constrained by commutation with the continuation structure. We show that composition at the survivor level is generically non-functorial, and that restoring compositional closure requires an extended ledger coordinate; in stable regimes this induces emergent time and a projectable causal order.

Building on this core, we formulate conditional realization criteria under which: (i)~stable refinement ceilings select low effective dimension; (ii)~order-dependent compression of phase-valued configurations reproduces interference patterns and Born-type weighting; and (iii)~exceptional projectability across boundaries implies holographic capacity bounds with area scaling. Computational evidence (Toy Model~5) shows that the dimensional selection threshold ($\tau = 4$) is first achieved at $n = 9$ on connected barbell graphs under locality-respecting extensions.

We further develop, as an explicit representational extension~(V7), a disciplined pathway to the Born rule. Non-functoriality of any representation attempt forces projectivization; projective non-contextual additivity then forces quadratic probability measures via a Gleason-class result. This derivation is strictly non-retroactive: it does not modify the foundational kernel.

Continuum structures enter only as asymptotic encodings of discrete invariants. No real numbers, manifolds, Hilbert spaces, or probability measures are assumed at the foundational layer. We conclude with a smuggling audit, a research program identifying open problems (including Einstein field equations and Standard Model gauge structure), and experimental discriminators.
\end{abstract}

\newpage

% ═══════════════════════════════════════════
% STATUS OF RESULTS
% ═══════════════════════════════════════════
\section*{Status of Results}
\addcontentsline{toc}{section}{Status of Results}

To prevent over-claiming, we distinguish four classes of results:

\textbf{Core Theorems (Unconditional):} Derived solely from Definitions~1.1--1.6 and finite set theory. These hold for any admissible extension class satisfying the axioms. Examples: Behavioral congruence equals isomorphism~(\ref{thm:behav-iso}), Unavoidability of compression~(\ref{thm:compress}), Entropy monotonicity~(\ref{thm:entropy}), Admissible law selection~(\ref{thm:law}), Non-functoriality~(\ref{thm:nonfunc}), Associative lift~(\ref{thm:lift}), Time emergence~(\ref{thm:time}).

\textbf{Conditional Realization Results:} Hold for specific classes of relational signatures satisfying stated structural criteria. Examples: Refinement ceiling selection~(3.11), Lorentzian signature~(3.21), Interference patterns~(4.5), Holographic bounds~(5.3).

\textbf{Representational Extension Results (V7):} Hold within the V7 representational layer, which introduces explicit new axioms (NF, RBA, PNA) above the core kernel. Non-retroactive: no V7 result modifies or strengthens any core or conditional result. Examples: Projectivization from non-functoriality~(V7-L1), Born-type rule~(V7-T1).

\textbf{Research Program (Open):} Conjectured results and identified derivation targets that have not yet been proved. Included to define the research frontier, not to claim completed work. Examples: Einstein field equations, Navier--Stokes structure, Standard Model gauge group, generation bounds.

\section*{Claim Legend}
\addcontentsline{toc}{section}{Claim Legend}
\begin{itemize}[leftmargin=1.5em]
  \item \claimtag{Unconditional} --- proved from the core kernel (Definitions~1.1--1.6) and finite set theory only.
  \item \claimtag{Conditional} --- proved only under additional structural hypotheses, all stated explicitly.
  \item \claimtag{Representational (V7)} --- proved only after adding NF/RBA/PNA; non-retroactive to lower layers.
  \item \claimtag{Open} --- conjectural target or research direction; not claimed as completed.
\end{itemize}

\newpage
\tableofcontents
\newpage

% ═══════════════════════════════════════════
% PART I
% ═══════════════════════════════════════════
\part{The Compression Framework}

\section{Primitive Ontology}

\begin{definition}[Relational Signature]\label{def:sig}
A \emph{relational signature} is a finite set $\mathcal{L} = \{R_i\}_{i \in I}$ where each $R_i$ has finite arity $k_i \in \mathbb{N}$.
\end{definition}

\begin{definition}[Micro-configuration]\label{def:micro}
A \emph{micro-configuration} is a finite $\mathcal{L}$-structure $\omega = (X, \{R_i^\omega\}_{i \in I})$ where $X$ is a finite set (the \emph{sites}) and $R_i^\omega \subseteq X^{k_i}$. Let $\mathbf{Str}_{\mathcal{L}}^{\mathrm{fin}}$ denote the class of all such structures.
\end{definition}

\begin{definition}[Configuration Space]\label{def:config}
A \emph{configuration space} is a finite subset $\Omega \subset \mathbf{Str}_{\mathcal{L}}^{\mathrm{fin}}$.
\end{definition}

\begin{definition}[Structural Equivalence]\label{def:equiv}
Two configurations $\omega, \omega' \in \Omega$ are \emph{structurally equivalent}, written $\omega \cong \omega'$, if there exists an isomorphism of $\mathcal{L}$-structures between them.
\end{definition}

\section{Admissible Extensions}

\begin{definition}[Extension Operator]\label{def:ext}
An \emph{extension operator} is a function $E \colon \Omega \to \mathbf{Str}_{\mathcal{L}_E}^{\mathrm{fin}}$ where $\mathcal{L}_E \supseteq \mathcal{L}$ is an expansion of the base signature.
\end{definition}

\begin{definition}[Admissible Extension Class]\label{def:admiss}
An \emph{admissible extension class} $\mathfrak{E}$ is a maximal class of extension operators satisfying:
\begin{enumerate}[label=(\roman*)]
  \item \textbf{Closure:} If $E_1, E_2 \in \mathfrak{E}$, then $E_2 \circ E_1 \in \mathfrak{E}$.
  \item \textbf{Congruence Preservation:} If $\omega \cong \omega'$, then $E(\omega) \cong E(\omega')$ for all $E \in \mathfrak{E}$.
  \item \textbf{Definability Closure:} If a partition $\mathcal{P}$ of $\Omega$ is invariant under all extensions in $\mathfrak{E}$, then $\mathcal{P}$ is \emph{incidence-definable}: there exists a first-order formula $\varphi(x)$ over the signature of $\mathfrak{A}$ such that $\omega, \omega'$ lie in the same cell of $\mathcal{P}$ iff $\mathfrak{A} \models \varphi(\omega) \leftrightarrow \varphi(\omega')$. No new predicates, second-order quantification, or external semantics are permitted; only the existing relational signature is used.
  \item \textbf{Effective Finiteness:} For all $E \in \mathfrak{E}$, the image $E(\Omega)$ has finitely many isomorphism types.
\end{enumerate}
\end{definition}

\begin{remark}
Condition~(iii) is the \emph{definability closure} axiom. It constrains the admissible class to be structurally complete: no extension-invariant partition of $\Omega$ can exist that is not already definable from the base $\mathcal{L}$-structure. This prevents artificial admissible classes, hidden partition injection, and entropy inflation. The formulation is purely syntactic (first-order definability over given relations) and does not introduce semantic primitives or model-theoretic ontology. Since any incidence-definable partition is automatically $\mathrm{Aut}(\mathfrak{A})$-invariant, definability closure aligns with automorphism reduction and prevents hidden gauge freedom. It does \emph{not}, by itself, force the induced symmetry group $G$ to act transitively on $\Sigma$ (see Remark~\ref{rem:transitivity} below).
\end{remark}

\begin{definition}[Locality-Respecting Admissible Class]\label{def:local-class}
Fix a graph metric $d_\Omega$ on each configuration's site set. An admissible extension class $\mathfrak{E}$ is \emph{$R$-locality-respecting} if there exists $R \in \mathbb{N}$ such that every primitive operation used by each $E \in \mathfrak{E}$ has support contained in a ball of radius $R$ in $d_\Omega$. Operations with unbounded support (for example, global complement) are excluded.
\end{definition}

\begin{remark}
Definition~\ref{def:local-class} provides the canonical formal meaning of ``R-local'' used later in Toy Model~5 and in the Part~VI diversity discussion.
\end{remark}

\begin{definition}[Behavioral Congruence]\label{def:behav}
Define a relation $\sim$ on $\Omega$ by:
\[
  \omega \sim \omega' \iff \forall E \in \mathfrak{E},\; E(\omega) \cong E(\omega').
\]
\end{definition}

\begin{lemma}\label{lem:equiv}
$\sim$ is an equivalence relation.
\end{lemma}
\begin{proof}
Reflexivity and symmetry are immediate. Transitivity follows from transitivity of $\cong$ and closure of $\mathfrak{E}$ under composition.
\end{proof}

\begin{theorem}[Behavioral Congruence Equals Isomorphism]\label{thm:behav-iso}
For any admissible extension class $\mathfrak{E}$:
\[
  \omega \sim \omega' \iff \omega \cong \omega'.
\]
That is, the behavioral congruence classes are exactly the isomorphism classes of $\mathcal{L}$-structures present in $\Omega$.
\end{theorem}
\begin{proof}
($\Leftarrow$) If $\omega \cong \omega'$, then $E(\omega) \cong E(\omega')$ for all $E \in \mathfrak{E}$ by congruence preservation (Definition~\ref{def:admiss}(ii)). Hence $\omega \sim \omega'$.

($\Rightarrow$) The identity operator $E_{\mathrm{id}}(\omega) := \omega$ satisfies all four conditions of Definition~\ref{def:admiss} and is therefore in every admissible class (it is also required by closure, since $E \circ E_{\mathrm{id}} = E$). If $\omega \sim \omega'$, then in particular $E_{\mathrm{id}}(\omega) \cong E_{\mathrm{id}}(\omega')$, i.e., $\omega \cong \omega'$.
\end{proof}
\assumptionfooter{Definitions~\ref{def:admiss}, \ref{def:behav}}{Kernel}{Proved.}

\begin{remark}\label{rem:survivor-iso}
Theorem~\ref{thm:behav-iso} implies that the survivor space $\Sigma$ is canonically identified with the set of isomorphism types realized in $\Omega$. Compression $q \colon \Omega \to \Sigma$ is the quotient by structural isomorphism. This holds for \emph{any} admissible extension class, not only for specific choices.
\end{remark}

\begin{definition}[Survivor Space]\label{def:survivor}
The \emph{survivor space} is the quotient $\Sigma := \Omega / {\sim}$. By Theorem~\ref{thm:behav-iso}, $\Sigma$ is the set of $\mathcal{L}$-isomorphism types present in $\Omega$. The quotient map $q \colon \Omega \to \Sigma$ is the \emph{compression map}.
\end{definition}

\begin{theorem}[Unavoidability of Compression]\label{thm:compress}
If $\Omega$ contains distinct but isomorphic configurations (i.e., distinct labeled representatives of the same isomorphism type), then $q$ is non-injective.
\end{theorem}
\begin{proof}
If $\omega \neq \omega'$ but $\omega \cong \omega'$, then $\omega \sim \omega'$ by Theorem~\ref{thm:behav-iso}, so $q(\omega) = q(\omega')$.
\end{proof}
\assumptionfooter{Theorem~\ref{thm:behav-iso}}{Kernel}{Proved.}

\section{Structural Entropy and Uniformity}

\begin{definition}[Structural Entropy]\label{def:entropy}
The \emph{structural entropy} of $\Omega$ is $S(\Omega) := \log |\Sigma|$. This is a combinatorial invariant measuring distinguishability capacity; any thermodynamic interpretation is regime-dependent.
\end{definition}

\begin{lemma}[Induced Quotient Map]\label{lem:quotient}
For any $E \in \mathfrak{E}$, the map $\bar{E} \colon \Sigma \to \Sigma_E$ defined by $\bar{E}([\omega]) := [E(\omega)]$ is well-defined and surjective.
\end{lemma}
\begin{proof}
If $[\omega] = [\omega']$, then $\omega \sim \omega'$, so $E(\omega) \sim_E E(\omega')$ by closure. Surjectivity follows from the definition of $\Sigma_E$.
\end{proof}

\begin{theorem}[Entropy Monotonicity]\label{thm:entropy}
For all admissible extensions $E \in \mathfrak{E}$:
\[
  S(E(\Omega)) \leq S(\Omega).
\]
\end{theorem}
\begin{proof}
By Lemma~\ref{lem:quotient}, $|\Sigma_E| \leq |\Sigma|$. Taking logarithms yields the result.
\end{proof}
\assumptionfooter{Definition~\ref{def:entropy}, Lemma~\ref{lem:quotient}}{Kernel}{Proved.}

\section{Admissible Laws and Structural Constraints}

\begin{definition}[Admissible Symmetry Group]\label{def:admiss-sym}
The \emph{admissible symmetry group} $G \leq \mathrm{Sym}(\Sigma)$ is the group of bijections of $\Sigma$ induced by admissible self-extensions that act bijectively on isomorphism types. The \emph{orbit decomposition} $\Sigma = O_1 \sqcup \cdots \sqcup O_m$ is the partition of $\Sigma$ into $G$-orbits.
\end{definition}

\begin{definition}[Admissible Law]\label{def:law}
An \emph{admissible law} is a bijection $\ell \colon \Sigma \to \Sigma$ that commutes with the admissible continuation structure: for every admissible extension $E \in \mathfrak{E}$ and every $\sigma \in \Sigma$,
\[
  \ell(\bar{E}(\sigma)) = \bar{E}(\ell(\sigma)),
\]
wherever both sides are defined.
\end{definition}

\begin{theorem}[Law Selection]\label{thm:law}
Assume distinct $G$-orbits are extension-separable: for every $i \neq j$, there exists $E \in \mathfrak{E}$ with $\bar{E}(O_i) \neq \bar{E}(O_j)$. Then every admissible law preserves each $G$-orbit setwise: if $\sigma \in O_i$, then $\ell(\sigma) \in O_i$. In particular, no admissible law exchanges distinct orbits.
\end{theorem}
\begin{proof}
An admissible law $\ell$ commutes with every element of $G$ (since $G$ is generated by admissible self-extensions). If $\ell(\sigma) \in O_j$ for some $\sigma \in O_i$, then for any $g \in G$ with $g\sigma = \sigma''$, commutativity gives $\ell(\sigma'') = g\ell(\sigma) \in O_j$, so $\ell(O_i) \subseteq O_j$. By bijectivity, $\ell(O_i)=O_j$. Apply commutation with an extension-separating $E$: if $\bar{E}(O_i) \neq \bar{E}(O_j)$, then $\ell \circ \bar{E} = \bar{E} \circ \ell$ is impossible on $O_i$. Contradiction. Therefore $\ell(O_i)=O_i$ for each $i$.
\end{proof}
\assumptionfooter{Definitions~\ref{def:admiss-sym}, \ref{def:law}; extension-separability hypothesis in Theorem~\ref{thm:law}}{Kernel}{Proved under stated separability condition.}

\begin{remark}[Structural exclusion is weight-free]\label{rem:weight-free}
Theorem~\ref{thm:law} excludes inadmissible laws on purely algebraic grounds: commutation with the continuation structure, not numerical weights or frequencies. No probability measure, stationary distribution, or counting argument is used. This is the primary structural constraint of the kernel.
\end{remark}

\begin{remark}[Non-transitivity and cross-orbit structure]\label{rem:transitivity}
The admissible symmetry group $G$ is not in general transitive on all of $\Sigma$. On 3-vertex simple graphs, for instance, $G \cong \mathbb{Z}_2$ (generated by complementation), with orbits $\{\text{empty}, \text{triangle}\}$ and $\{\text{single-edge}, \text{path}\}$ (Appendix~\ref{app:toy-models}, Toy Model~1). The orbit structure is determined by the continuation structure of $\mathfrak{E}$. No admissible law can exchange orbits; this is a structural fact, not a probabilistic one.
\end{remark}

\subsection{Admissible Weightings (Diagnostic)}

The following results characterize admissible probability assignments on $\Sigma$. These are \emph{diagnostic}: they describe what weightings are consistent with admissibility, but no exclusion of laws or symmetries depends on them.

\begin{definition}[Admissible Weighting]\label{def:weight}
A function $w \colon \Sigma \to \mathbb{R}_{\geq 0}$ is an \emph{admissible weighting} if (i)~$\sum_{\sigma \in \Sigma} w(\sigma) = 1$, and (ii)~$w$ is invariant under all bijections of $\Sigma$ induced by admissible self-extensions.
\end{definition}

\begin{theorem}[Forced Uniformity Within Orbits]\label{thm:uniform}
Any admissible weighting $w$ is uniform within each $G$-orbit: if $\sigma, \sigma' \in O_i$, then $w(\sigma) = w(\sigma')$.
\end{theorem}
\begin{proof}
$G$ acts transitively within each orbit (by definition). Invariance of $w$ under $G$ gives $w(\sigma) = w(g\sigma) = w(\sigma')$ for any $g \in G$ with $g\sigma = \sigma'$.
\end{proof}
\assumptionfooter{Definition~\ref{def:weight}, group action of $G$ on each orbit}{Kernel}{Proved.}

\begin{corollary}[Structure of Admissible Weightings]\label{cor:weight}
Every admissible weighting has the form
\[
  w(\sigma) = \frac{\lambda_i}{|O_i|} \quad \text{for } \sigma \in O_i,
\]
where $\lambda_1, \dots, \lambda_m \geq 0$ satisfy $\sum_i \lambda_i = 1$. The parameters $\lambda_i$ are cross-orbit weights (empirical input); the intra-orbit uniformity $1/|O_i|$ is forced by admissibility.
\end{corollary}

\begin{remark}[Decision rule]\label{rem:decision}
No exclusion of laws, symmetries, or identities in this monograph depends on numerical weights unless those weights are uniquely forced by admissibility. After Theorem~\ref{thm:law}, all structural exclusions are algebraic.
\end{remark}


% ═══════════════════════════════════════════
% PART II
% ═══════════════════════════════════════════
\part{Non-Functoriality and the Emergence of Time}

\section{Induced Maps and Typed Composition}

\begin{definition}[Induced Survivor Maps]\label{def:induced}
For an admissible extension $E \in \mathfrak{E}$, define the induced map $\bar{E} \colon \Sigma \to \Sigma_E$ by $\bar{E}([\omega]) := [E(\omega)]$. For composable extensions $E_1, E_2$, denote by $\bar{E}_2^{(E_1)} \colon \Sigma_{E_1} \to \Sigma_{E_2 \circ E_1}$ the induced map of $E_2$ acting on the $E_1$-extended configuration space.
\end{definition}

\begin{theorem}[Non-Functoriality]\label{thm:nonfunc}
In general:
\[
  \overline{E_2 \circ E_1} \neq \bar{E}_2^{(E_1)} \circ \bar{E}_1
\]
as maps $\Sigma \to \Sigma_{E_2 \circ E_1}$.
\end{theorem}
\begin{proof}
See Appendix~\ref{app:toy-models} (Toy Model~1).
\end{proof}

\section{Defect and the Ledger}

\begin{definition}[Defect Indicator]\label{def:defect}
For composable $E_1, E_2 \in \mathfrak{E}$ and $\sigma \in \Sigma$, define
\[
  \delta(E_1, E_2; \sigma) :=
  \begin{cases}
    0 & \text{if } \overline{E_2 \circ E_1}(\sigma) = (\bar{E}_2^{(E_1)} \circ \bar{E}_1)(\sigma), \\
    1 & \text{otherwise.}
  \end{cases}
\]
\end{definition}

\begin{definition}[Ledger Coordinate]\label{def:ledger}
For a composable history $\mathbf{E} = (E_1, \dots, E_k)$, the \emph{cumulative defect} starting at $\sigma \in \Sigma$ is
\[
  \Delta(\mathbf{E}; \sigma) := \sum_{i=1}^{k-1} \delta(E_i, E_{i+1}; \sigma_i),
\]
where $\sigma_1 := \sigma$ and $\sigma_{i+1} := \overline{E_i \circ \cdots \circ E_1}(\sigma)$. The \emph{ledger coordinate} $n \in \mathbb{N}$ records this cumulative obstruction count.
\end{definition}

\begin{definition}[Lifted Survivor Space]\label{def:lifted}
The \emph{lifted survivor space} is $\widetilde{\Sigma} := \Sigma \times \mathbb{N}$. For a history $\mathbf{E}$, the lifted action is
\[
  \widetilde{\Phi}(\mathbf{E})(\sigma, n) := \bigl(\bar{\mathbf{E}}(\sigma),\; n + \Delta(\mathbf{E}; \sigma)\bigr).
\]
\end{definition}

\begin{theorem}[Associative Lift]\label{thm:lift}
For composable histories $\mathbf{E}, \mathbf{F}$: the lifted composite equals the composite of lifts, provided
\[
  \Delta(\mathbf{F} \circ \mathbf{E}; \sigma) = \Delta(\mathbf{E}; \sigma) + \Delta(\mathbf{F}; \bar{\mathbf{E}}(\sigma)).
\]
\end{theorem}
\begin{proof}
Direct computation. The survivor component agrees by definition of composition. The ledger component agrees by the additivity assumption on $\Delta$.
\end{proof}

\section{Dynamical Regimes and Time Emergence}

\subsection{Histories as Finite Admissible Words}

We eliminate temporal interpretation from the term ``history.''

\begin{definition}[Admissible Word]\label{def:word}
An \emph{admissible word} on a region $U$ is a finite sequence $w = (E_1, \dots, E_n)$ with each $E_i \in \mathfrak{E}$ supported in $U$. Its action is the composite extension $E_w := E_n \circ \cdots \circ E_1$. The \emph{length} of $w$ is $|w| := n$.
\end{definition}

\begin{remark}
Words are purely algebraic objects in the semigroup generated by $\mathfrak{E}$. No temporal interpretation is required. All phase definitions, obstruction arguments, and ceiling results refer only to finite words and their defect densities. When we write ``history'' below, it is shorthand for ``admissible word.''
\end{remark}

\subsection{Structural Phase Definitions}

\begin{definition}[Structural Phases]\label{def:phases}
Phases are defined purely by certified structural properties of regions, not by probabilistic or dynamical assumptions. Fix a region $U$, canonical probe family $\mathcal{P}_{R_0}$, and locality radius $R_0$.
\begin{itemize}[nosep]
  \item \textbf{Phase~A} (Certified Stability): $U$ admits LS-2 certification (Definition~\ref{def:LS2}), recurrence certification with coverage epoch $T_{\mathrm{cov}}(U)$ (Definition~\ref{def:reccert}), small certified defect floor $D_{\mathrm{floor}}(U) \ll 1$ (Definition~\ref{def:floor}), and floor-attaining admissible words with small slack $\xi \leq \xi_{\max} \ll 1$ (Definition~\ref{def:slack}).
  \item \textbf{Phase~B} (Obstruction-Dominated): LS-2 and recurrence certification hold, but $D_{\mathrm{floor}}(U)$ is not small or floor-attaining slack cannot be kept small. Refinement persistence is unstable.
  \item \textbf{Phase~C} (Non-Certifiable): At least one of LS-2 certification, recurrence certification, or existence of floor-attaining words fails. No dimensional-selection mechanism is justified.
\end{itemize}
\end{definition}

\begin{remark}
The phase classification depends only on certified combinatorial properties---obstruction witness counts, probe coverage bounds, and defect densities of finite admissible words. No probability measure, temperature parameter, or ergodic principle is used. The earlier pre-adaptation and mild-selection hypotheses (Assumptions~\ref{ass:preadapt},~\ref{ass:mild}) provide one sufficient mechanism by which Phase-A conditions may be realized; they are not logically required for the phase definitions or for $\tau$-selection.
\end{remark}

\begin{theorem}[Time Emergence]\label{thm:time}
Local clocks---monotone functions of the ledger coordinate $n$---exist only in Phase~A regimes.
\end{theorem}
\begin{proof}
In Phase~B, $\Delta n$ fluctuates; in Phase~C, $\Delta n$ is non-monotone due to reorganizing transitions. Only Phase~A admits stable monotone accumulation.
\end{proof}

\begin{definition}[Causal Order]\label{def:causal}
For survivor classes $\sigma, \sigma' \in \Sigma$: $\sigma \prec \sigma'$ iff there exists an admissible evolution modifying $\sigma$ that changes $\sigma'$.
\end{definition}

\begin{theorem}[Emergent Causality]\label{thm:causal}
In Phase~A, $\prec$ defines a stable partial order.
\end{theorem}
\begin{proof}
Bounded influence ensures antisymmetry and transitivity under admissibility.
\end{proof}


% ═══════════════════════════════════════════
% PART III
% ═══════════════════════════════════════════
\part{Geometric Emergence (Conditional Realization)}

\statusnotice{The results in this part are conditional: they hold only in Phase-A regimes admitting bounded propagation, locally separating admissible extensions, and low dynamical defect density. The central claim is not that higher-dimensional refinements are combinatorially impossible, but that they are dynamically unstable under Phase-A evolution.}

\section{Refinements and Projectability}

\begin{definition}[Refinement]\label{def:refine}
A \emph{refinement} is a map $r \colon \Omega \to \Omega_r$ that augments configurations with additional relational structure, increasing distinguishability at the micro-configuration level.
\end{definition}

\begin{definition}[Projectable Refinement]\label{def:project}
A refinement $r$ is \emph{projectable} if for every admissible extension $E \in \mathfrak{E}$:
\[
  \overline{E \circ r} = \bar{E} \circ \bar{r}.
\]
Projectable refinements survive long-term compression and admit effective representation at the survivor level.
\end{definition}

\begin{definition}[Independent Refinements]\label{def:indep}
A collection $\{r_1, \dots, r_k\}$ of projectable refinements is \emph{independent} if no $r_i$ factors through the joint refinement of the others.
\end{definition}

\begin{definition}[Refinement Ceiling]\label{def:ceiling}
The \emph{refinement ceiling} $\tau$ is the maximal cardinality of a set of mutually independent, projectable refinements admitted by a given Phase-A regime.
\end{definition}

\begin{theorem}[Finiteness of the Refinement Ceiling]\label{thm:finite}
For any finite configuration space $\Omega$, the refinement ceiling $\tau$ is finite.
\end{theorem}
\begin{proof}
Each projectable refinement strictly increases survivor distinguishability. By entropy monotonicity (Theorem~\ref{thm:entropy}), survivor cardinality is bounded. Hence refinement chains must terminate.
\end{proof}

\section{Local Load and Defect Generation}

The obstruction to unlimited refinement is not global entropy, but \emph{local non-functoriality}. More independent refinements create more local distinctions, which create more opportunities for order-dependent obstruction.

\begin{definition}[Local Load]\label{def:load}
Let $U$ be a radius-$R_0$ neighborhood in a Phase-A region. The \emph{local load} $L(U)$ is the number of distinct joint refinement labels realized within $U$.
\end{definition}

\begin{definition}[Local Probe Family]\label{def:probe}
A \emph{canonical local probe family} $\mathcal{P}_{R_0}$ is a finite collection of admissible extensions supported in radius-$R_0$ neighborhoods, each adding definable tags determined by local isomorphism type (degree, neighbor-degree multiset, short-cycle counts, WL-1 color classes). All probes are invariant under $\mathrm{Aut}(\mathfrak{A})$ and compatible with definability closure.
\end{definition}

\begin{assumption}[LS-0: Local Probe Completeness]\label{ass:LS0}
For the chosen locality scale $R_0$, the admissible class $\mathfrak{E}$ contains a canonical local probe family $\mathcal{P}_{R_0}$ with the property: if two lifted states are locally distinguishable (non-isomorphic on some radius-$R_0$ ball), then some probe $P \in \mathcal{P}_{R_0}$ distinguishes them.
\end{assumption}

\begin{definition}[LS-2: Pairwise Obstruction Sufficiency]\label{def:LS2}
A region $U$ satisfies \emph{pairwise obstruction sufficiency} if for any two distinct refinement labels $a \neq b$ realized in $U$, there exist probes $P_a, P_b \in \mathcal{P}_{R_0}$ and a survivor class $\sigma$ supported in $U$ such that applying the probes in different orders yields a non-functorial obstruction: $\delta(P_a, P_b; \sigma) = 1$ or $\delta(P_b, P_a; \sigma) = 1$.
\end{definition}

\begin{remark}
LS-2 is strictly weaker than full local separability. It asserts only that distinct local labels generate at least one order-dependent obstruction witness, which is exactly what the obstruction-growth argument requires. By Lemma~\ref{lem:LS2-suff} below, LS-2 holds whenever distinct labels produce mutually interfering probes with differential collapse timing---the mechanism exhibited by Toy Model~1 and verified on all tested substrates.
\end{remark}

\begin{lemma}[Sufficient Criterion for LS-2]\label{lem:LS2-suff}
Fix a canonical local probe family $\mathcal{P}_{R_0}$. If for any two distinct local labels $a \neq b$ there exist probes $P_a, P_b \in \mathcal{P}_{R_0}$ and a survivor class $\sigma$ such that (i)~both probes act nontrivially on $\sigma$, (ii)~applying one alters the local-type conditions tested by the other, and (iii)~the induced survivor maps collapse at different stages under the two orders, then LS-2 holds.
\end{lemma}

\begin{theorem}[Local Obstruction Growth (Linear Form)]\label{thm:obstruct}
Let $U$ be a Phase-A region with local load $L(U) = L$. Assume pairwise obstruction sufficiency LS-2 holds on $U$. Then there exists a constant $c_0 \in (0,1]$ depending only on $\mathcal{P}_{R_0}$ such that every probe-rich admissible history $\mathbf{E}$ supported in $U$ contains at least $c_0(L-1)$ obstruction events.
\end{theorem}
\begin{proof}[Proof sketch]
Fix a distinguished label $a_0$ realized in $U$. For each of the $L-1$ other labels $a \neq a_0$, LS-2 provides a probe pair whose relative order produces an obstruction witness on some survivor class supported in $U$. Each requirement contributes at least one obstruction event. The constant $c_0$ accounts for probe reuse and overlaps.
\end{proof}

\begin{remark}\label{rem:growth}
The linear bound $\Omega(L)$ replaces the earlier heuristic $f(L) = \binom{L}{2}/(\binom{L}{2}+1)$. It is structurally safer: no genericity or density argument is required, only the finitary LS-2 condition. For $g(k) = 2^k$ branching, even linear growth suffices to force $\tau \leq 4$.
\end{remark}


\section{Certified Obstruction Floors and Recurrence Bounds}

This section constructs a purely structural defect floor from LS-2 certificates and finite recurrence bounds. No dynamical selection principle, probability measure, or pre-adaptation hypothesis is required.

\begin{definition}[Certified Obstruction Count]\label{def:ncert}
Let $U$ be a region satisfying LS-2 (Definition~\ref{def:LS2}) with canonical probe family $\mathcal{P}_{R_0}$. The \emph{certified obstruction count} $N_{\mathrm{cert}}(U)$ is the number of distinct label pairs $(a,b)$ for which an explicit obstruction witness (probe pair + survivor class) has been verified. By Theorem~\ref{thm:obstruct}, $N_{\mathrm{cert}}(U) \geq c_0(L(U) - 1)$.
\end{definition}

\begin{definition}[Recurrence Certificate]\label{def:reccert}
A region $U$ admits a \emph{recurrence certificate} with bound $T_{\mathrm{cov}}(U)$ if: for every admissible word $w$ supported in $U$ with $|w| \geq T_{\mathrm{cov}}(U)$, every probe $P \in \mathcal{P}_{R_0}$ occurs at least once in $w$. The bound $T_{\mathrm{cov}}(U)$ is the minimal such length. This is a finite combinatorial property of the admissible class restricted to $U$.
\end{definition}

\begin{definition}[Certified Defect Floor]\label{def:floor}
For a region $U$ with certified obstruction count $N_{\mathrm{cert}}(U)$ and coverage epoch $T_{\mathrm{cov}}(U)$, the \emph{certified defect floor} is
\[
  D_{\mathrm{floor}}(U) := \frac{c_\star \cdot N_{\mathrm{cert}}(U)}{T_{\mathrm{cov}}(U)},
\]
where $c_\star > 0$ is a structural constant depending only on the probe family. This is the minimum defect density that any probe-covering admissible word on $U$ must exhibit.
\end{definition}

\begin{definition}[Floor-Attaining Slack]\label{def:slack}
For an admissible word $w$ supported in $U$, the \emph{floor-attaining slack} is
\[
  \xi(w; U) := \overline{\mathcal{D}}(w) - D_{\mathrm{floor}}(U).
\]
By definition, $\xi \geq 0$. A word is \emph{floor-attaining} if $\xi(w; U) \leq \xi_{\max}$ for some small $\xi_{\max} \ll 1$.
\end{definition}

\begin{remark}
The certified floor is a structural lower bound on defect density, not an assumption about realized dynamics. It follows from finite combinatorics: any admissible word long enough to cover all probes must encounter the certified obstruction witnesses. The floor replaces the earlier defect budget $D_{\max}$ as the load-bearing inequality in $\tau$-selection.
\end{remark}


\section{Phase-A Mechanisms and Refinement Ceiling Selection}

Structural Pre-Adaptation and Mild Defect Aversion provide one mechanism by which realized admissible words remain floor-attaining (Definition~\ref{def:slack}). They are not required for the obstruction lower bound, the certified defect floor, or the $\tau$-selection inequality itself---which compares structural obstruction pressure against the certified floor plus realized slack.

\begin{definition}[Defect Accumulation Functional]\label{def:defectfunc}
For a history $\mathbf{E} = (E_1, \dots, E_n)$ acting on survivor class $\sigma$, the \emph{normalized defect density} is
\[
  \overline{\mathcal{D}}(\mathbf{E}; \sigma) := \frac{1}{n} \sum_{i=1}^{n-1} \delta(E_i, E_{i+1}; \sigma_i).
\]
\end{definition}

\begin{definition}[Local Extension Neighborhood]\label{def:neigh}
Fix a locality-respecting admissible class $\mathfrak{E}$. For a lifted state $\tilde{\sigma}$, let $\mathcal{N}(\tilde{\sigma}) := \{ E \in \mathfrak{E} \mid E \text{ is supported within a bounded-radius neighborhood of } \tilde{\sigma} \}$. By effective finiteness, $\mathcal{N}(\tilde{\sigma})$ is finite.
\end{definition}

\begin{definition}[Step Defect Score]\label{def:stepdefect}
Let $\mathcal{R}(\tilde{\sigma})$ denote the set of projectable refinements available at $\tilde{\sigma}$. For $E \in \mathcal{N}(\tilde{\sigma})$, the normalized step defect is
\[
  d(E; \tilde{\sigma}) := \frac{\#\{ r \in \mathcal{R}(\tilde{\sigma}) \mid r \text{ is destroyed by } E \}}{|\mathcal{R}(\tilde{\sigma})|}.
\]
This is a rational number determined purely by finite counting.
\end{definition}

\begin{definition}[Low-Defect Fraction and Preservation Basin]\label{def:basin}
For $\varepsilon \in [0,1]$, define the low-defect fraction $\alpha_\varepsilon(\tilde{\sigma}) := |\{ E \in \mathcal{N}(\tilde{\sigma}) \mid d(E; \tilde{\sigma}) \leq \varepsilon \}| / |\mathcal{N}(\tilde{\sigma})|$. Fix parameters $\varepsilon \ll 1$ and $\eta \ll 1$. A set of lifted states $\mathcal{B} \subseteq \widetilde{\Sigma}$ is a \emph{preservation basin} if:
\begin{enumerate}[label=(\roman*)]
  \item For all $\tilde{\sigma} \in \mathcal{B}$, $\alpha_\varepsilon(\tilde{\sigma}) \geq 1 - \eta$.
  \item \emph{Trap asymmetry}: if $E \in \mathcal{N}(\tilde{\sigma})$ satisfies $d(E; \tilde{\sigma}) > \varepsilon$, then the successor state $\tilde{\sigma}'$ satisfies $\alpha_\varepsilon(\tilde{\sigma}') \leq \eta$.
\end{enumerate}
\end{definition}

\begin{assumption}[Structural Pre-Adaptation]\label{ass:preadapt}
There exist parameters $\varepsilon \ll 1$, $\eta \ll 1$ and a nonempty preservation basin $\mathcal{B}$.
\end{assumption}

\begin{remark}
Assumption~\ref{ass:preadapt} is purely structural. It asserts no dynamics, probabilities, or weights---only combinatorial geometry of local admissible neighborhoods. Toy Model~6 verifies this property on barbell and connected substrates: low-defect moves form ${\sim}79\%$ of local neighborhoods in tested basins.
\end{remark}

\begin{assumption}[Mild Defect Aversion]\label{ass:mild}
There exists an effective evolution rule $\Phi$ selecting realized histories such that, whenever $E, E' \in \mathcal{N}(\tilde{\sigma})$ satisfy $d(E; \tilde{\sigma}) < d(E'; \tilde{\sigma})$, the rule $\Phi$ does not systematically prefer $E'$ over $E$.
\end{assumption}

\begin{remark}
Assumption~\ref{ass:mild} introduces no probability measure, temperature parameter, or numerical weighting. It only excludes systematic preference for higher-defect steps when lower-defect steps are available.
\end{remark}

\begin{assumption}[No Probe Starvation]\label{ass:nps}
In Phase-A preservation basins, the realized evolution $\Phi$ does not permanently suppress any probe in $\mathcal{P}_{R_0}$ on stable regions $U$: for each probe $P \in \mathcal{P}_{R_0}$ and each such $U$, there exist arbitrarily long realized histories in which $P$ occurs at least once.
\end{assumption}

\begin{theorem}[Emergent Defect Budget]\label{thm:budget}
Under Structural Pre-Adaptation (Assumption~\ref{ass:preadapt}) and Mild Defect Aversion (Assumption~\ref{ass:mild}), any realized history $\mathbf{E} = (E_1, \dots, E_n)$ remaining entirely within the preservation basin $\mathcal{B}$ satisfies $\overline{\mathcal{D}}(\mathbf{E}) \leq \varepsilon + O(\eta)$.
\end{theorem}
\begin{proof}
Within $\mathcal{B}$, at least a $(1-\eta)$ fraction of $\mathcal{N}(\tilde{\sigma})$ consists of $\varepsilon$-defect steps. By Mild Defect Aversion, higher-defect steps are not systematically preferred. Since high-defect steps lead to states with $\alpha_\varepsilon \leq \eta$ (trap asymmetry), histories remaining in $\mathcal{B}$ must overwhelmingly consist of $\varepsilon$-defect steps. Averaging over $n$ steps yields the bound.
\end{proof}

\begin{corollary}[Derived Phase-A Budget]\label{cor:budget}
In a preservation basin, the defect-budget inequality $\overline{\mathcal{D}} \leq D_{\max}$ holds with $D_{\max} := \varepsilon + O(\eta)$. The budget is a \emph{derived regime property}, not a primitive axiom.
\end{corollary}

\begin{definition}[Phase-A Failure Gates]\label{def:failgates}
Phase-A stability fails if any of the following occur:
\begin{enumerate}[label=(F\arabic*)]
  \item \textbf{No Pre-Adaptation:} $\alpha_\varepsilon(\tilde{\sigma})$ is bounded away from $1$ for all candidate stable states and all small $\varepsilon$.
  \item \textbf{No Trap Asymmetry:} High-defect steps do not reduce $\alpha_\varepsilon$ in successor states.
  \item \textbf{Adversarial Selection:} Realized histories systematically prefer higher-defect steps.
  \item \textbf{Probe Starvation:} Some canonical probe is permanently suppressed, preventing obstruction-witness generation.
\end{enumerate}
\end{definition}

\begin{definition}[Local Branching Function]\label{def:branch}
Let $g(k)$ denote the minimal local load induced by $k$ independent projectable refinements within radius-$R_0$ neighborhoods:
\[
  L(U) \geq g(k).
\]
Under a local product condition, $g(k) = 2^k$; in general, only monotonic growth is assumed.
\end{definition}

\begin{remark}\label{rem:product}
Under a local product condition---if for each subset $S \subseteq \{1, \dots, k\}$ there exists a radius-2 neighborhood realizing a configuration where the joint label differs exactly on $S$---then $L \geq 2^k$. This is a checkable structural hypothesis, not a universal claim.
\end{remark}

\begin{theorem}[Refinement Ceiling Selection (Certified Form)]\label{thm:ceiling}
Let $U$ be a Phase-A region (Definition~\ref{def:phases}) satisfying:
\begin{enumerate}[label=(\roman*)]
  \item bounded propagation with radius $R_0$ (derived from Phase-A; Theorem~\ref{thm:locality});
  \item LS-2 pairwise obstruction sufficiency on $U$ (Definition~\ref{def:LS2}), yielding certified obstruction count $N_{\mathrm{cert}}(U)$;
  \item recurrence certification with bound $T_{\mathrm{cov}}(U)$ (Definition~\ref{def:reccert}), yielding certified defect floor $D_{\mathrm{floor}}(U)$ (Definition~\ref{def:floor});
  \item floor-attaining admissible words with slack $\xi \leq \xi_{\max}$ (Definition~\ref{def:slack});
  \item local branching $L(U) \geq g(k)$ for $k$ independent refinements.
\end{enumerate}
Then any $k$ satisfying
\[
  \frac{c_0(g(k) - 1)}{|w|} > D_{\mathrm{floor}}(U) + \xi_{\max}
\]
for floor-attaining admissible words $w$ of characteristic length cannot persist on $U$. In particular, if $g(5)$ is large enough that the obstruction pressure from $5$ refinements exceeds the certified floor plus slack, then $\tau(U) \leq 4$.
\end{theorem}
\begin{proof}
By recurrence certification, any admissible word of length $|w| \geq T_{\mathrm{cov}}(U)$ covers all canonical probes. By Theorem~\ref{thm:obstruct}, $k$ independent refinements induce local load $L \geq g(k)$, generating at least $c_0(g(k)-1)$ obstruction events in any probe-covering word. Dividing by word length yields a lower bound on defect density. For a floor-attaining word with slack $\xi \leq \xi_{\max}$, the realized defect density satisfies $\overline{\mathcal{D}}(w) \leq D_{\mathrm{floor}}(U) + \xi_{\max}$. If the obstruction lower bound exceeds this, the refinements cannot persist. Under typical local product conditions, $g(5) \geq 32$, giving $c_0 \cdot 31$ obstruction events---sufficient to exceed $D_{\mathrm{floor}} + \xi_{\max}$ on standard epoch scales.
\end{proof}

\begin{remark}[Phase-A stability selection]\label{rem:selection}
The ceiling $\tau \leq 4$ is a Phase-A stability selection: exceeding it forces obstruction density above the certified floor plus attainable slack. Four independent refinements correspond to three reversible (spatial) directions plus one irreversible ledger direction; a fifth drives obstruction pressure beyond any small certified floor. The $\tau$-selection inequality involves no dynamical primitives: it compares a structural obstruction lower bound against a structural defect floor, both computed from finite combinatorics.
\end{remark}

\begin{definition}[Effective Dimension]\label{def:dim}
The \emph{effective dimension} $d$ is the maximal number of independent projectable spatial refinements.
\end{definition}

\begin{corollary}[$3+1$ Emergence]\label{cor:31}
Under the conditions of Theorem~\ref{thm:ceiling}, $d = 3$ spatial dimensions and $1$ temporal dimension emerge from $\tau = 4$.
\end{corollary}

\section{Emergent Lorentzian Signature}

The distinction between reversible spatial refinements and the irreversible ledger direction induces a metric signature. We formalize this as a representation-theoretic result on invariant quadratic forms.

\begin{definition}[Local Flow Space]\label{def:flow}
Assume a Phase-A region $U \subseteq \widetilde{\Sigma}$ admits a coarse-grained encoding $\Xi \colon U \to \mathbb{R}^{d+1}$ such that the image of small-step histories from $\tilde{\sigma}$ defines a pointed cone of increments $C_{\tilde{\sigma}}$. Let $V_{\tilde{\sigma}} = \operatorname{span}(C_{\tilde{\sigma}})$ be the \emph{local flow space}.
\end{definition}

\begin{definition}[Temporal/Spatial Decomposition]\label{def:decomp}
Let $\tau_{\tilde{\sigma}} \in V_{\tilde{\sigma}}^*$ be a \emph{clock form} satisfying $\tau_{\tilde{\sigma}}(v) > 0$ for all $v \in C_{\tilde{\sigma}} \setminus \{0\}$, derived from the ledger projection. Define the \emph{temporal line} $T_{\tilde{\sigma}}$ as the extremal direction in $C$ with respect to $\tau$, and the \emph{spatial subspace} $S_{\tilde{\sigma}} := \ker(\tau_{\tilde{\sigma}}) \cap V_{\tilde{\sigma}}$. Then $V_{\tilde{\sigma}} = T_{\tilde{\sigma}} \oplus S_{\tilde{\sigma}}$.
\end{definition}

\begin{assumption}[Reversibility Structure]\label{ass:reverse}
The realized near-symmetry group $G_{\tilde{\sigma}} \leq \mathrm{GL}(V_{\tilde{\sigma}})$ satisfies:
\begin{enumerate}[label=(\roman*)]
  \item $G_{\tilde{\sigma}}$ acts transitively on the unit sphere of $S_{\tilde{\sigma}}$ (spatial isotropy);
  \item $G_{\tilde{\sigma}}$ fixes $T_{\tilde{\sigma}}$ setwise and preserves $C_{\tilde{\sigma}}$.
\end{enumerate}
\end{assumption}

\begin{definition}[Admissibly Invariant Quadratic Form]\label{def:quadform}
A quadratic form $Q_{\tilde{\sigma}}$ on $V_{\tilde{\sigma}}$ is \emph{admissibly invariant} if $Q_{\tilde{\sigma}}(gv) = Q_{\tilde{\sigma}}(v)$ for all $g \in G_{\tilde{\sigma}}$, $v \in V_{\tilde{\sigma}}$.
\end{definition}

\begin{lemma}[Spatial Positivity]\label{lem:spatial}
If $Q$ is admissibly invariant and nondegenerate on $S$, then $Q|_S$ is definite.
\end{lemma}
\begin{proof}
Spatial isotropy (transitivity on spatial directions) forces $Q$ to have constant sign on $S \setminus \{0\}$. If $Q$ were indefinite, transitivity would carry a direction with $Q > 0$ to one with $Q < 0$, contradicting invariance.
\end{proof}

\begin{lemma}[Block Diagonality]\label{lem:block}
Under the same hypotheses, $Q(t, s) = 0$ for all $t \in T$, $s \in S$.
\end{lemma}
\begin{proof}
For any $g \in G$, $gt = \pm t$ (setwise preservation of $T$) and $gs$ ranges over all directions in $S$ (transitivity). Invariance gives $Q(t, s) = Q(gt, gs) = \pm Q(t, gs)$. As $gs$ ranges over $S$, the linear functional $s \mapsto Q(t, s)$ must be constant on the sphere; the only such linear functional is zero.
\end{proof}

\begin{lemma}[Temporal Sign]\label{lem:temporal}
If $Q$ is nondegenerate, block-diagonal with respect to $T \oplus S$, positive on $S \setminus \{0\}$, and $\partial C \subseteq \{v : Q(v) = 0\}$ (cone--null compatibility), then $Q$ is negative on $T \setminus \{0\}$.
\end{lemma}
\begin{proof}
If $Q$ were positive on $T$, then $Q$ would be positive definite on $V$, with null set $= \{0\}$. But cone--null compatibility requires a nontrivial null boundary (the boundary of $C$ is nonempty). Contradiction.
\end{proof}

\begin{theorem}[Lorentzian Signature]\label{thm:lorentz}
Let $U$ be a Phase-A region with the decomposition $V = T \oplus S$, spatial isotropy (Assumption~\ref{ass:reverse}), and an admissibly invariant nondegenerate quadratic form $Q$ satisfying cone--null compatibility ($\partial C \subseteq \{v : Q(v) = 0\}$). Then, up to overall scale, $Q$ has Lorentzian signature $(-,+,+,+)$, with the negative direction aligned with $T$.
\end{theorem}
\begin{proof}
Lemma~\ref{lem:spatial} fixes spatial positivity. Lemma~\ref{lem:block} gives block diagonality. Lemma~\ref{lem:temporal} forces the temporal sign negative. Together: one negative eigenvalue (temporal), $d$ positive eigenvalues (spatial). For $d = 3$, this is $(-,+,+,+)$.
\end{proof}

\begin{remark}
The negative sign is not postulated; it is forced by (i)~a distinguished irreversible monotone direction, (ii)~isotropy on reversible directions, and (iii)~an invariant causal cone represented by a quadratic form.
\end{remark}

\section{Emergent Locality}

Locality is not assumed as a kernel axiom. We show that Phase-A preservation basins structurally exclude support-unbounded (global) admissible moves.

\begin{definition}[Support Radius]\label{def:supp}
For an extension $E$ acting on a configuration with carrier $X$, the \emph{support radius} $\mathrm{rad}(E)$ is the smallest $R$ such that $E$ is determined by induced substructures on radius-$R$ balls and modifies relations only within those balls. If no such $R$ exists, set $\mathrm{rad}(E) = \infty$. Call $E$ \emph{global} if $\mathrm{rad}(E) = \infty$.
\end{definition}

\begin{definition}[Global Sensitivity]\label{def:sens}
A state $\tilde{\sigma}$ is \emph{globally sensitive} if there exists $c > 0$ such that every global move $E \in \mathcal{N}(\tilde{\sigma})$ with $\mathrm{rad}(E) = \infty$ satisfies $d(E; \tilde{\sigma}) \geq c$.
\end{definition}

\begin{theorem}[Phase-A Implies Effective Locality]\label{thm:locality}
Assume Structural Pre-Adaptation (Assumption~\ref{ass:preadapt}) with parameters $\varepsilon \ll 1$, $\eta \ll 1$, and suppose states in the preservation basin $\mathcal{B}$ are globally sensitive with constant $c > \varepsilon$. If global moves occur in a non-negligible fraction $\gamma$ of the one-step neighborhood (i.e., $|\mathcal{N}_{\mathrm{glob}}(\tilde{\sigma})|/|\mathcal{N}(\tilde{\sigma})| \geq \gamma$), then necessarily $\gamma \leq \eta$. In particular, for $\eta \ll 1$, a preservation basin cannot coexist with a one-step neighborhood containing a non-negligible fraction of global moves.
\end{theorem}
\begin{proof}
Let $\tilde{\sigma} \in \mathcal{B}$. By basin definition, $\alpha_\varepsilon(\tilde{\sigma}) \geq 1 - \eta$. By global sensitivity with $c > \varepsilon$, every global move has defect $> \varepsilon$ and therefore does not contribute to the low-defect numerator. Hence $\alpha_\varepsilon(\tilde{\sigma}) \leq 1 - \gamma$. Combining: $1 - \eta \leq 1 - \gamma$, i.e., $\gamma \leq \eta$.
\end{proof}

\begin{corollary}[Locality as Basin Necessity]\label{cor:locality}
If $\eta \ll 1$ in a Phase-A preservation basin, then either (i)~global moves are absent from $\mathcal{N}(\tilde{\sigma})$, or (ii)~global moves form a vanishing fraction. In either case, Phase-A stability selects effectively local admissible dynamics.
\end{corollary}

\begin{remark}[No converse claimed]
Theorem~\ref{thm:locality} does not assert that locality implies Phase-A stability. Local admissible classes may still fail to admit preservation basins if refinements are sparse or trap asymmetry is absent. Locality is therefore a necessary, but not sufficient, condition for Phase-A.
\end{remark}


% ═══════════════════════════════════════════
% PART IV
% ═══════════════════════════════════════════
\part{Quantum Mechanics Without Quantization}

\statusnotice{Sections~\ref{sec:meas}--\ref{sec:interf} hold for configuration spaces admitting phase-valued relational structures. Section~\ref{sec:born} is a representational extension~(V7) with explicitly introduced axioms.}

\section{Measurement as Forced Projection}\label{sec:meas}

\begin{definition}[Measurement Context]\label{def:meas}
A \emph{measurement context} is a refinement attempt $r$ that exceeds the refinement ceiling $\tau$.
\end{definition}

\begin{theorem}[Measurement as Forced Projection]\label{thm:meas}
Any refinement exceeding $\tau$ is erased under admissible evolution, projecting onto the survivor space $\Sigma$.
\end{theorem}
\begin{proof}
By definition of $\tau$, such refinements cannot be projectable. Compression forces projection.
\end{proof}

\section{Interference and Probability}\label{sec:interf}

\begin{construction}[Phase-Valued Configurations]\label{con:phase}
Consider micro-configurations $\omega = (\gamma, \phi, x)$ where $\gamma \in \{L, R\}$ (path label), $\phi \in \mathbb{Z}_M$ (structural phase encoded in auxiliary cyclic relational attachments), and $x \in X$ (position).
\end{construction}

\begin{theorem}[Structural Interference]\label{thm:interf}
When which-path information is erased, the multiplicity at position $x$ is
\[
  |q^{-1}(\sigma_x)| = 2M\bigl(1 + \cos\theta(x)\bigr),
\]
where $\theta(x) = 2\pi I(x)/M$ and $I(x)$ is the structural phase offset.
\end{theorem}
\begin{proof}
Phase differences modulo $M$ create behavioral congruence classes. The count of compatible phase pairs varies as the discrete Fourier transform of $\mathbb{Z}_M$, yielding cosine modulation.
\end{proof}

\begin{theorem}[Probability from Compression]\label{thm:prob}
The probability weighting compatible with admissibility is
\[
  P(\sigma) = \frac{|q^{-1}(\sigma)|}{|\Omega|}.
\]
\end{theorem}
\begin{proof}
Micro-configurations are distinguishable prior to compression. Admissibility forbids distinguishing survivors beyond multiplicity. Relative frequencies are forced.
\end{proof}

\section{The Born Rule via Representational Extension (V7)}\label{sec:born}

\statusnotice{This section introduces the V7 representational extension. It adds three explicit axioms (NF, RBA, PNA) above the core kernel. No V7 result modifies any core or conditional result. The firewall is strict and non-retroactive.}

\subsection{Non-Functoriality Axiom (NF)}

The core kernel establishes non-functoriality as a structural fact (Theorem~\ref{thm:nonfunc}). V7 elevates this to a representational principle.

\begin{axiom}[NF: Non-Functoriality]\label{ax:nf}
Any representation attempt $F$ from the admissibility/witness domain $\mathcal{W}$ into a linear representation domain $\mathcal{R}$ exhibits an irreducible obstruction $\Omega$ such that:
\begin{enumerate}[label=(\roman*)]
  \item For some composable morphisms $a, b$ in $\mathcal{W}$, $F(b \circ a) \neq F(b) \circ F(a)$.
  \item The deviation is not eliminable by admissible redefinition or post-composition with an isomorphism.
  \item $\Omega$ is stable under admissible extension.
\end{enumerate}
\end{axiom}

\subsection{Representation Bridge Axiom (RBA)}

\begin{axiom}[RBA: Representation Bridge]\label{ax:rba}
There exists a representation domain $\mathcal{H}$ with:
\begin{enumerate}[label=(\roman*)]
  \item Objects are finite-dimensional complex inner-product spaces.
  \item Morphisms include linear maps up to phase-compatible equivalence.
  \item A probability functional $P$ can be defined on rays (projective classes) in each object.
\end{enumerate}
\end{axiom}

\noindent\textbf{Anti-Ontology Clause.} RBA licenses a representational codomain only. Nothing in RBA asserts that Hilbert-space structure is ontologically fundamental.

\subsection{Projectivization from Non-Functoriality}

\begin{lemma}[NF Forces Projectivization (V7-L1)]\label{lem:project}
Given NF and RBA, any admissible representation that remains stable under admissible extension must factor through the projective space of the representation domain. That is, meaning cannot be assigned to vectors, only to rays.
\end{lemma}
\begin{proof}
Under NF, different compositional paths yield images differing by the obstruction $\Omega$. PASS forbids resolving this by gauge choice or preferred paths. The maximal quotient removing the obstruction without PASS violation is the projective quotient. Only ray-level structure survives.
\end{proof}

\subsection{Projective Non-Contextual Additivity (PNA)}

\begin{axiom}[PNA: Projective Non-Contextual Additivity]\label{ax:pna}
There exists a function $p \colon \mathrm{Rays}(H) \to [0,1]$ such that:
\begin{enumerate}[label=(\roman*)]
  \item For every admissible context $\{[e_i]\}$ (orthonormal basis representing a maximal compatible witness decomposition in $\mathcal{W}$), $\sum_i p([e_i]) = 1$.
  \item $p([v])$ depends only on the ray $[v]$, not on which context contains it.
  \item $p$ is invariant under admissible extension (PASS compatible).
\end{enumerate}
Non-contextuality expresses invariance under admissible refactorizations of the same underlying witness decomposition.
\end{axiom}

\subsection{Born-Type Rule}

\begin{theorem}[Born-Type Probability (V7-T1)]\label{thm:born}
Assume NF (Axiom~\ref{ax:nf}), RBA (Axiom~\ref{ax:rba}), PNA (Axiom~\ref{ax:pna}), and $\dim(H) \geq 3$. Then any admissible probability assignment on rays must take the form
\[
  p([v]) = \frac{\langle v | \rho | v \rangle}{\langle v | v \rangle}
\]
for a unique positive semidefinite operator $\rho$ with $\mathrm{Tr}(\rho) = 1$. For pure states $\rho = |\psi\rangle\langle\psi| / \langle\psi|\psi\rangle$, this reduces to the standard Born rule:
\[
  p([v]) = \frac{|\langle v | \psi \rangle|^2}{\langle v | v \rangle \langle \psi | \psi \rangle}.
\]
\end{theorem}
\begin{proof}
By Lemma~\ref{lem:project}, probability must be defined on rays. PNA imposes non-contextual additivity across orthogonal decompositions. By a Gleason-class result (a mathematical theorem about measures on projective spaces, not a physical assumption), the only such function for $\dim \geq 3$ is quadratic in the ray representative.
\end{proof}
\assumptionfooter{Axioms~\ref{ax:nf}, \ref{ax:rba}, \ref{ax:pna}; Lemma~\ref{lem:project}; Gleason-class theorem for $\dim \ge 3$}{Representational (V7)}{Proved conditional on V7 axioms.}

\begin{remark}[Dimension caveat]
For $\dim(H) = 2$, additional conditions are required. This is a known structural fact about projective measures, not a defect of the framework.
\end{remark}

\begin{remark}[Scope boundary for V7]
Section~\ref{sec:born} derives only the measure form on rays under NF+RBA+PNA. It does \emph{not} derive Hamiltonian dynamics, decoherence mechanisms, measurement-update dynamics, or a preferred ontology for Hilbert structure.
\end{remark}


% ═══════════════════════════════════════════
% PART V
% ═══════════════════════════════════════════
\part{Holography and Cosmological Emergence}

\statusnotice{The results in this section hold for boundary-bulk systems satisfying exceptional projectability.}

\section{Exceptional Projectability}

\begin{definition}[Boundary-Bulk System]\label{def:bulk}
A system $(\Omega_\partial, \Omega_{\mathrm{bulk}}, I_h)$ with \emph{holographic incidence} $I_h \subseteq X_\partial \times X_{\mathrm{bulk}}$.
\end{definition}

\begin{definition}[Exceptional Projectability]\label{def:except}
The system is \emph{exceptionally projectable} if every bulk refinement factors through a boundary refinement.
\end{definition}

\begin{theorem}[Holographic Bound]\label{thm:holo}
In exceptionally projectable regimes:
\[
  S_{\mathrm{bulk}} \leq \frac{\tau}{2}\, S_\partial.
\]
\end{theorem}
\begin{proof}
Bulk survivor classes are determined by boundary projections. Each boundary class supports at most $\tau/2$ independent bulk refinements (ceiling division). Taking logs yields the entropy bound.
\end{proof}

\begin{corollary}[Area Law --- Model-Dependent]\label{cor:area}
When $S_\partial \propto A$ (boundary area), we obtain $S_{\mathrm{bulk}} \leq A/4G$ with $1/4G \sim (\tau/2) \cdot c$, where $c$ is a realization-dependent constant.
\end{corollary}

\section{Cosmological Emergence}

\begin{definition}[Emergence Boundary]\label{def:emerge}
$n_\star := \min\{n : \text{the lifted evolution at step } n \text{ enters the spacetime regime}\}$.
\end{definition}

\begin{theorem}[Cosmological Emergence]\label{thm:cosmo}
Spacetime begins at $n_\star$. Prior to $n_\star$, the system is in Phase~B or~C, lacking stable monotone ledger accumulation.
\end{theorem}
\begin{proof}
Geometric definability requires Phase-A stability, which is absent prior to $n_\star$.
\end{proof}

\begin{corollary}[Big Bang Reinterpretation]\label{cor:bang}
The Big Bang is not a singular instant but the onset of geometric definability.
\end{corollary}


% ═══════════════════════════════════════════
% PART VI: RESEARCH PROGRAM
% ═══════════════════════════════════════════
\part{Research Program (Open Problems)}

The following are identified derivation targets and conjectured results. They define the research frontier of NFC. None are claimed as completed.

\section{Strengthening Dimensional Selection}

\textbf{Status:} Conditional theorem (Theorem~\ref{thm:ceiling}), with $\tau = 4$ now achieved computationally on the barbell $K_3\text{-}b_3\text{-}K_3$ ($n = 9$) under locality-respecting radius-$2$ extensions in the sense of Definition~\ref{def:local-class} (Toy Model~5).

\textbf{Key finding (Toy Model~5):} The refinement ceiling $\tau$ scales linearly with bridge length for the barbell family, reaching $\tau = 4$ at $n = 9$ and growing without bound. Under the unrestricted maximal class, $\tau = 0$; locality is necessary and sufficient for non-trivial $\tau$.

\textbf{Remaining open problems:}
\begin{enumerate}[nosep]
  \item \textbf{Minimum degree condition.} The local load $g \geq 32$ is met at clique-adjacent vertices but not at degree-2 bridge midpoints ($g_{\min} = 30$ on $K_4$-lad$_2$-$K_4$). Formalize the minimum-degree requirement as a condition on the local separability hypothesis.
  \item \textbf{Locality axiomatization.} Adopt Definition~\ref{def:local-class} as the canonical locality condition and instantiate it at radius $R=2$ for current computations. The complement boundary suggests: admissible extensions act within bounded-radius neighborhoods.
  \item \textbf{Defect analysis on barbells.} Compute defect densities for $\tau = 4, 5, 6$ on the barbell family.
  \item \textbf{Full $g \geq 32$ closure.} Determine whether degree-$\geq 3$ ladder-bridge structures with multi-arity refinements achieve $g_{\min} \geq 32$ uniformly.
\end{enumerate}

\textbf{Priority:} High. The $\tau = 4$ achievement provides the first direct computational support for the dimensional selection mechanism on a connected substrate.

\subsection*{V7.5 Integration: Dimension-Selection Enumeration Protocol}\label{sec:v75-enum}
To integrate the V7.5 dimension-selection program into the canonical monograph, we fix an explicit finite enumeration protocol for claims about threshold attainment and scaling. This subsection is methodological: it does not introduce new kernel axioms and does not upgrade any open claim to theorem status.

\textbf{Protocol inputs.}
\begin{enumerate}[label=\textbf{E\arabic*.},nosep]
  \item A finite graph family $\mathcal{F}_n$ at each size $n$ (exhaustive for small $n$, curated families for larger $n$).
  \item A fixed locality-respecting admissible class $\mathfrak{E}_{R}$ as in Definition~\ref{def:local-class} with radius parameter $R$ (currently $R=2$ baseline).
  \item A refinement alphabet specification (vertex-only, then multi-arity tiers).
  \item A canonical survivor-independence criterion (GF(2) rank over surviving projectable refinements).
\end{enumerate}

\textbf{Enumeration outputs.}
\begin{enumerate}[label=\textbf{O\arabic*.},nosep]
  \item $\tau_{\max}(n;\mathfrak{E}_R,\mathcal{A})$: maximal certified refinement ceiling at size $n$.
  \item Witness family achieving $\tau_{\max}$, with orbit profile and local-load profile.
  \item Boundary-collapse diagnostics under global augmentations (e.g., complement) to certify locality sensitivity.
  \item Confidence class: exhaustive / certificate-backed / heuristic (must be explicitly tagged in tables).
\end{enumerate}

\textbf{Current integrated status (V7.5 baseline).} Under radius-2 locality-respecting extensions, the barbell family provides certified threshold attainment at $n=9$ with $\tau=4$, and linear bridge-length scaling thereafter (Toy Model~5). Under global augmentation classes, collapse to $\tau=0$ persists as the controlling boundary effect.

\textbf{Reproducibility hooks (repository scripts).} The V7.5 pipeline is implemented by the local analysis scripts `multi_arity_extraction.py`, `q5_q12_diagnosis.py`, `ladder_certificate_v2.py`, `lc_explicit.py`, and `verify_preadapt_ls2.py`. These scripts define the operational meaning of the enumeration/certificate steps and should be treated as the canonical computational companion for this subsection.

\textbf{Governance rule.} Any future statement of the form ``minimum $n$ for $\tau\ge k$'' must cite: (i) extension class, (ii) refinement alphabet tier, (iii) witness family, and (iv) confidence class. Statements missing one of these are downgraded to exploratory notes.

\section{Maximality, Transitivity, and Local Diversity}\label{sec:diversity}

\textbf{Status:} Partially resolved. Transitivity settled (Theorem~\ref{thm:behav-iso}, Theorem~\ref{thm:law}). Local diversity characterized computationally with barbell threshold result (Toy Model~5); theoretical questions remain open.

\textbf{Resolved.} Behavioral congruence equals isomorphism in every admissible class. Admissible laws must preserve $G$-orbits on purely algebraic grounds (Theorem~\ref{thm:law}). The admissible symmetry group $G$ is generically non-transitive (Toy Model~1: $G \cong \mathbb{Z}_2$ on 3-vertex graphs). Cross-orbit weighting is empirical input; all structural exclusions are weight-free (Remark~\ref{rem:weight-free}).

\textbf{New finding: barbell threshold and scaling.} Toy Model~5 reveals a sharp hierarchy: $\tau = 0$ under global operations but $\tau$ grows without bound under locality-respecting radius-$2$ extensions (Definition~\ref{def:local-class}). The barbell family $K_3$-$b_N$-$K_3$ achieves $\tau = 4$ at $n = 9$ and $\tau = 6$ at $n = 13$, confirming linear scaling. Complement is the critical boundary: adding it collapses $\tau$ to~0 at all sizes.

\textbf{Candidate strengthening (revised).}

\begin{conjecture}[Diversity-Forcing Maximality]\label{conj:diversity}
For a locality-respecting admissible extension class $\mathfrak{E}$ (extensions supported within bounded-radius neighborhoods) on sufficiently large configurations ($n \geq n_0$), the number of inequivalent surviving refinement types on a patch of size $k$ is at least $\Omega(\log k)$ (weak form) or $\Omega(k^\varepsilon)$ for some fixed $\varepsilon > 0$ (strong form).
\end{conjecture}

The barbell data provides partial support: $\tau$ grows linearly with $n$ on the barbell family, consistent with the weak form. The conjecture is explicitly conditioned on locality; without it, $\tau = 0$.

\textbf{Test vectors:} Test edge-coloring refinements on barbells. Compute local load under multi-color refinements. Determine whether the $g(\tau)$ gap closes with richer refinement types.

\textbf{Priority:} High. This question is upstream of all conditional realization results.

\section{Einstein Field Equations}

\textbf{Status:} Not yet attempted within the consolidated framework.

\textbf{Conjectured form:} Under Phase-A coherence, uniform-compression extremum, and minimal geometric encoding, the unique tensor-form obstruction balance is isomorphic to $G_{\mu\nu} = \kappa\, T_{\mu\nu}$.

\textbf{Required:} Precise structural inputs analogous to Lovelock-type conditions. A comparison identifying which assumptions are kernel-derived and which are empirically selected. A gate provenance audit for all encoding assumptions.

\section{Navier--Stokes Structure}

\textbf{Status:} Not yet attempted within the consolidated framework.

\textbf{Conjectured form:} Under bilinear flux forcing and cycle-defect dissipation, the unique refinement-stable transport balance reproduces Navier--Stokes.

\section{Standard Model Gauge Structure}

\textbf{Status:} Identified as open problem. Not yet attempted as formal derivation.

\textbf{Conjectured form:} $\mathrm{SU}(3) \times \mathrm{SU}(2) \times \mathrm{U}(1)$ appears as the unique continuous compression summary of a finite stabilizer action on a defect-typed regime alphabet.

\textbf{Required:} A precise argument showing why these specific Lie groups emerge. Gate provenance audit distinguishing kernel-derived from physically-encoded inputs.

\section{Generation Bound}

\textbf{Status:} Not yet derived.

\textbf{Required:} Independent derivation of the witness capacity from the kernel.

\section{Computational Complexity of Behavioral Congruence}

\textbf{Status:} Open. GI-completeness conjectured.

\subsection*{Complexity Roadmap}
\begin{enumerate}[label=\textbf{C\arabic*.}]
  \item \textbf{Reduction baseline.} Build explicit many-one reductions between behavioral congruence on restricted signatures and graph isomorphism instances.
  \item \textbf{Parameterized regimes.} Classify fixed-parameter behavior under bounded arity, bounded degree, and bounded locality radius.
  \item \textbf{Benchmark families.} Maintain canonical finite families (cycles, ladders, barbells, multi-arity ladders) with exact instance generators and expected complexity profiles.
  \item \textbf{Certificate formats.} Develop machine-checkable witness/counterexample certificates aligned with LS-2 and recurrence certification.
\end{enumerate}

\appendix
\section{Toy Models: Empirical Evidence for Finite Refinement Ceiling and Stability Selection}
\label{app:toy-models}

This appendix consolidates four small finite toy models constructed entirely within the core kernel of NFC (Definitions~\ref{def:sig}--\ref{def:project}, finite set theory only; no conditional regimes, Phase distinctions, or V7 representational extensions are invoked). All computations are exact and manual. The models provide concrete, reproducible evidence bearing on two central structural results:

\begin{enumerate}
  \item The admissible symmetry group $G$ is generically non-transitive on $\Sigma$ (Toy Model~1), confirming Theorem~\ref{thm:behav-iso} and the orbit structure described in Remark~\ref{rem:transitivity}.
  \item Projectable refinements are finite and are strongly selected toward globally uniform or low-complexity structures under natural admissible extensions, providing evidence for finite refinement ceilings and stability selection toward low effective dimension (Theorem~\ref{thm:finite}).
\end{enumerate}

\subsection{Toy Model 1: 3-Vertex Graphs --- Non-Transitive Survivor Partition}

\textbf{Signature:} $\mathcal{L}=\{R\}$ (binary symmetric irreflexive; undirected simple graphs).\\
\textbf{Carrier:} 3 vertices.\\
\textbf{Micro-configurations:} $\omega_\triangle$ (K$_3$), $\omega_\vee$ (path of length 2).

\textbf{Admissible class} $\mathfrak{E}_0$ (candidate maximal, closed under composition): $E_{\mathrm{id}}$, $E_{\mathrm{neigh}}$ (neighborhood-type unaries), $E_{\mathrm{deg}}$ (degree unaries), $E_{\mathrm{comp}}$ (component-size unaries), $E_{\mathrm{symdiff}}$ (complement binary), $E_{\mathrm{univ}}$ (universal binary), $E_{\mathrm{trans}}$ (transitive closure binary), $E_{\mathrm{clique}}$ (maximal-clique-size unaries), $E_{\mathrm{aut}}$ (automorphism orbit-size unaries).

\textbf{Results:}
\begin{itemize}[nosep]
  \item Behavioral congruence $\sim$ yields \textbf{two survivor classes}: $[\triangle]$ (high-symmetry triangle-like) and $[\vee]$ (asymmetric path-like).
  \item No composition forces $\omega_\triangle \sim \omega_\vee$.
  \item The induced symmetry monoid on $\Sigma$ contains collapsing maps but no invertible map that swaps the two classes.
  \item $G$ does \textbf{not} act transitively.
  \item The partition survives aggressive coarse-graining and appears forced by base incidence (connectivity/transitivity differences), consistent with the intent of Definition~\ref{def:admiss}(iii).
\end{itemize}

\textbf{Kernel-level statement:} The non-transitivity of $G$ is not a failure of maximality. By Theorem~\ref{thm:behav-iso}, behavioral congruence equals isomorphism in \emph{every} admissible class, so non-isomorphic structures can never be merged. The partition $\{[\triangle], [\vee]\}$ is forced by base incidence and is consistent with condition~(iii). The admissible symmetry group $G$ acts transitively within each orbit but not across orbits; cross-orbit weighting is empirical input (Remark~\ref{rem:transitivity}).

\textbf{Non-functoriality} (sampled composable pairs): present in approximately 40--50\% of tested cases (primarily due to historical label preservation and differential collapse timing).

\subsection{Toy Model 2: 4-Vertex Graphs --- Unary Refinement Erasure}

\textbf{Carrier:} 4 vertices. \textbf{Signature:} $\mathcal{L}=\{R\}$.\\
\textbf{Base graphs:} $C_4$ (cycle), $K_4$ (complete), $P_4$ (path).\\
\textbf{Refinements} (unary predicate $C$):
\begin{itemize}[nosep]
  \item $r_1$: uniform (all vertices $C$)
  \item $r_2$: balanced bipartite (alternating $C/\neg C$)
  \item $r_3$: unbalanced 3-valued ($1{+}1{+}2$ pattern)
\end{itemize}

\textbf{Class:} $\mathfrak{E}_0$ extended with $E_{\mathrm{labelmerge}}$ (merge labels on identical neighborhoods) and $E_{\mathrm{labelsym}}$ (merge labels with identical neighbor multisets \emph{and} orbit sizes).

\textbf{Results:}
\begin{itemize}[nosep]
  \item $r_1$ (uniform): fully projectable; commutes with all operators; survives indefinitely.
  \item $r_2$ (bipartite): survives original $\mathfrak{E}_0$ but is \textbf{erased} under $E_{\mathrm{labelsym}} \circ E_{\mathrm{trans}}$ (neighbor multisets equalize on complete background, forcing merge to uniform).
  \item $r_3$ (3-valued): reduced to two labels by $E_{\mathrm{labelmerge}}$, then erased under $E_{\mathrm{labelsym}}$ + connectivity.
\end{itemize}

\textbf{Kernel-level statement:} Non-trivial unary refinements (local symmetry breaking, unbalanced multiplicity) are erased under label-behavioral coarse-graining combined with connectivity. Only globally uniform labeling remains stable. This gives direct evidence that the refinement ceiling $\tau$ is finite and small ($\tau \leq 2$; likely $\tau=1$ under maximal symmetry enforcement).

\subsection{Toy Model 3: 4-Vertex Graphs with Binary Relation --- Directed Structure Collapse}

\textbf{Signature:} $\mathcal{L}=\{R,S\}$ ($R$ undirected simple, $S$ binary not necessarily symmetric).\\
\textbf{Refinements on $S$:}
\begin{itemize}[nosep]
  \item $r_4$: directed 4-cycle
  \item $r_5$: two disjoint directed edges (matching)
\end{itemize}

\textbf{Class:} previous extensions + $E_{\mathrm{dirsym}}$ (add reverse arrows to $S$, merge equivalent post-symmetrization patterns).

\textbf{Results:}
\begin{itemize}[nosep]
  \item Both $r_4$ and $r_5$ \textbf{collapse immediately} under $E_{\mathrm{dirsym}}$ (directionality erased; $S$ reduces to undirected).
  \item Connectivity accelerates erasure.
  \item Directed binary refinements contribute 0 to $\tau$.
\end{itemize}

\textbf{Kernel-level statement:} Directed and asymmetric binary structure is more fragile than unary local breaking. Symmetrization dominates, favoring undirected, globally uniform patterns.

\subsection{Toy Model 4: 4-Vertex Graphs with Cyclic/Phase Attachments --- Periodic Structure Collapse}

\textbf{Signature:} $\mathcal{L}=\{R\}$ (extended with phase-valued unary predicate $\Phi$: values in $\mathbb{Z}_4 = \{0,1,2,3\}$).\\
\textbf{Base graphs:} $C_4$ (cycle), $K_4$ (complete).\\
\textbf{Refinements} (phase patterns on vertices):
\begin{itemize}[nosep]
  \item $r_6$: uniform phase ($\Phi=0$ everywhere)
  \item $r_7$: sequential cyclic ($0$-$1$-$2$-$3$ around $C_4$)
  \item $r_8$: period-2 repeating ($0$-$1$-$0$-$1$)
  \item $r_9$: unbalanced cluster ($0$-$0$-$1$-$2$)
\end{itemize}

\textbf{Class:} previous extensions + $E_{\mathrm{phasefold}}$ (merge phases equivalent under rotation or reflection if graph symmetry permits) and $E_{\mathrm{labelsym}}$ (adapted to phase values).

\textbf{Results:}
\begin{itemize}[nosep]
  \item $r_6$ (uniform): fully projectable; survives indefinitely (equivalent to uniform unary).
  \item $r_7$--$r_9$ (non-trivial cyclic/periodic): start with 2--4 distinct phases/neighborhood types, but \textbf{collapse} under connectivity ($E_{\mathrm{trans}}$) + symmetry merge ($E_{\mathrm{labelsym}}$ or $E_{\mathrm{phasefold}}$); neighbor multisets equalize on complete background, forcing full merge to uniform.
  \item Final effective diversity: \textbf{1} in all non-trivial cases.
\end{itemize}

\textbf{Kernel-level statement:} Cyclic and periodic phase structures are at least as fragile as directed binary refinements. Periodicity and local variation are erased by symmetry enforcement and connectivity. Non-trivial phase patterns contribute 0 to $\tau$.


\subsection{Toy Model 5: Refinement Ceiling Scaling and the Barbell Threshold}\label{sec:tm5}

\textbf{Purpose:} Systematic computation of the actual refinement ceiling $\tau$ on $n$-vertex simple graphs ($n = 3$--$13$) under extension classes of varying strength, testing whether the dimensional selection mechanism (Theorem~\ref{thm:ceiling}) can engage at accessible scales.

\textbf{Method.} For each graph, we: (i)~compute the automorphism orbit structure via color refinement, (ii)~enumerate all non-trivial projectable 2-color vertex refinements (colorings constant on orbits), (iii)~test survival under each extension, and (iv)~count independent survivors via GF(2) rank. Extension classes of increasing strength are tested:
\begin{itemize}[nosep]
  \item \textbf{R=1 local:} identity, degree-based edge filtering, pendant removal, local complementation.
  \item \textbf{Radius-2 locality-respecting:} Radius-1 operations plus transitive closure, common-neighbor edge addition, 2-core extraction (Definition~\ref{def:local-class}).
  \item \textbf{R=2 + complement:} R=2 plus graph complement and compositions.
  \item \textbf{Nuclear (maximal):} all of the above plus delete-all, complete-all.
\end{itemize}

\textbf{Phase 1: Small graphs ($n \leq 6$, exhaustive).}

\begin{center}
\begin{tabular}{lccccc}
\toprule
\textbf{Extension class} & \textbf{$n{=}3$} & \textbf{$n{=}4$} & \textbf{$n{=}5$} & \textbf{$n{=}6$} & \textbf{Best graph ($n{=}6$)} \\
\midrule
R=1 local         & 0 & 2 & 3 & 4 & $(2,2,2,2,3,3)$ \\
Radius-2 local      & 0 & 0 & $2^*$ & 3 & $(1,1,3,3,3,3)$ \\
R=2 + complement & 0 & 0 & 0 & 0 & --- \\
Nuclear          & 0 & 0 & 0 & 0 & --- \\
\bottomrule
\end{tabular}
\end{center}

{\footnotesize ${}^*$Achieved only on disconnected graphs ($K_3 \sqcup K_2$); first \emph{connected} radius-2 locality-respecting survivor at $n = 6$ (Definition~\ref{def:local-class}).}

\textbf{Phase 2: The barbell family (connected, $n = 7$--$13$).} Targeted search over graph families reveals that barbell graphs---two cliques connected by a bridge path---achieve the highest $\tau$ under locality-respecting radius-$2$ extensions (Definition~\ref{def:local-class}). The key discovery:

\begin{center}
\begin{tabular}{lcccc}
\toprule
\textbf{Graph} & \textbf{$n$} & \textbf{Orbits} & $\boldsymbol{\tau}$ \textbf{(R=2)} & $\boldsymbol{g(\tau)}$ \\
\midrule
$K_3\text{-}b_1\text{-}K_3$ & 7 & 3 & 2 & 2 \\
$K_3\text{-}b_2\text{-}K_3$ & 8 & 2 & 2 & 2 \\
$K_3\text{-}b_3\text{-}K_3$ & 9 & 4 & \textbf{4} & 3 \\
$K_3\text{-}b_4\text{-}K_3$ & 10 & 4 & \textbf{4} & 3 \\
$K_3\text{-}b_5\text{-}K_3$ & 11 & 5 & \textbf{5} & 3 \\
$K_3\text{-}b_7\text{-}K_3$ & 13 & 6 & \textbf{6} & 3 \\
$K_3\text{-}b_4\text{-}K_4$ & 11 & 8 & \textbf{5} & 3 \\
\bottomrule
\end{tabular}
\end{center}

\textbf{$\tau = 4$ is achieved at $n = 9$ on the connected barbell $K_3\text{-}b_3\text{-}K_3$.} This is the minimal connected structure meeting the threshold required by Theorem~\ref{thm:ceiling}. Furthermore, $\tau$ scales linearly with bridge length: $\tau \approx \lceil b/2 \rceil + 1$ for $K_3$-$b_N$-$K_3$.

\textbf{Why barbells work.} The barbell structure resists radius-2 locality-respecting extensions because:
(a)~cliques are internally stable (transitive closure is a no-op on complete subgraphs);
(b)~bridge vertices at different distances from the cliques are in distinct orbits with distinct neighborhoods;
(c)~degree-based filtering separates bridge (degree 2) from clique boundary vertices (degree $\geq 3$), preserving rather than erasing orbit distinctions.

\textbf{Local load gap: multi-arity analysis.} Under 2-color vertex refinements, $g(\tau) = 3$, far below $g(5) \geq 32$. Systematic exploration of richer refinement types reveals a structured path toward closing this gap:

\begin{center}
\begin{tabular}{lcc}
\toprule
\textbf{Refinement type} & $\boldsymbol{g_{\min}}$ & $\boldsymbol{g_{\max}}$ \\
\midrule
2-color vertex only & 3 & 4 \\
Combined vertex + edge (2-color each) & 5 & 8 \\
Multi-arity (V+E+P2+P3+T+C4), $K_3$-$b_N$-$K_3$ & 7 & 16 \\
Multi-arity, ladder bridge $K_4$-lad$_2$-$K_4$ & \textbf{30} & 34 \\
Multi-arity, ladder bridge $K_4$-lad$_4$-$K_4$ & 21 & 46 \\
Multi-arity, $K_4$-lad$_4$-$K_4$, R=3 neighborhoods & 23 & --- \\
\bottomrule
\end{tabular}
\end{center}

The bottleneck is always the lowest-degree vertex (bridge midpoint, degree~2), whose R=2 ball contains only~5 vertices and can support at most~${\sim}20$ distinct orbit types regardless of refinement arity. Ladder bridges (degree~3 throughout) raise $g_{\min}$ to~30. The clique-adjacent vertices already exceed $g = 32$: e.g., $K_4$-lad$_2$-$K_4$ achieves $g_{\max} = 34$, and $K_3$-lad$_5$-$K_5$ achieves $g_{\max} = 56$.

The remaining gap is \emph{structural and honest}: the local separability hypothesis in Theorem~\ref{thm:ceiling} implicitly requires minimum vertex degree $\geq 3$. This is a natural condition---it asserts that every vertex participates in at least three independent relational channels.

\textbf{Findings:}
\begin{enumerate}[nosep]
  \item \textbf{Locality is the critical variable.} Under radius-2 locality plus complement or nuclear classes, $\tau = 0$ at all sizes. Under radius-2 locality-respecting classes, $\tau$ grows without bound. Complement (a global operation) is the primary erasure mechanism.
  \item \textbf{The $\tau = 4$ threshold is achievable.} Barbell $K_3\text{-}b_3\text{-}K_3$ at $n = 9$ achieves $\tau = 4$ under locality-respecting radius-$2$ extensions (Definition~\ref{def:local-class}). This confirms that the dimensional selection mechanism can engage on modest connected structures.
  \item \textbf{Unbounded scaling.} $\tau$ grows linearly with bridge length for the barbell family, confirming that locality-respecting extension classes support arbitrarily high refinement ceilings.
  \item \textbf{Local load closes under multi-arity refinements on ladder bridges.} With vertex, edge, path, triangle, and 4-cycle orbit types, $g_{\min} = 30$ on $K_4$-lad$_2$-$K_4$ ($n = 12$). The $g \geq 32$ threshold is met at clique-adjacent vertices; the remaining gap at bridge midpoints requires minimum degree $\geq 3$, a natural structural condition.
\end{enumerate}

\textbf{Kernel-level statement:} Bounded propagation is a structural prerequisite for dimensional selection. Under the unrestricted maximal class, $\tau = 0$; under locality-respecting classes, $\tau$ grows without bound and reaches the Theorem~\ref{thm:ceiling} threshold ($\tau = 4$) at $n = 9$. The barbell $K_3\text{-}b_3\text{-}K_3$ is the minimal connected realization substrate.


\subsection{Toy Model 6: Defect Selection Principle Emergence}\label{app:toymodel6}

\statusnotice{Computational experiment ($n = 9$, $300$ Monte Carlo trials). Tests whether the Defect Selection Principle (DSP)---the assumption that Phase-A dynamics preferentially follow low-defect histories---emerges from the geometry of the configuration space rather than requiring independent postulation.}

\textbf{Setup.} Substrate: $K_3\text{-}b_3\text{-}K_3$ ($n = 9$, $\tau = 4$ under radius-2 locality-respecting extensions, $14$ non-trivial projectable refinements). At each time step, a random radius-2 locality-respecting extension (single-edge modification within a radius-2 ball) is applied. The \emph{defect} of a step is the fraction of refinements broken by the extension; the \emph{cumulative defect} of a history is the time-average of per-step defects. We compare three dynamics: (i)~uniform random; (ii)~defect-minimizing (greedy); (iii)~Boltzmann-weighted with inverse temperature $\beta$.

\textbf{Finding 1: Landscape pre-adaptation.} Of $28$ distinct radius-2 locality-respecting extensions from $K_3\text{-}b_3\text{-}K_3$, \textbf{22 (79\%) are defect-free}. Extensions that do break refinements exhibit all-or-nothing behavior: they break exactly $8/14$ refinements simultaneously, or none. The configuration space is geometrically pre-adapted for low-defect dynamics.

\textbf{Finding 2: Bimodal attractor structure.} Under uniform random dynamics ($300$ trials, $15$ steps), the final state is \emph{bimodal}: $42\%$ of trials preserve all $14$ refinements (full-preservation basin), $22\%$ destroy all refinements (full-destruction basin), and $36\%$ are partial. Mean cumulative defect: $0.35$. The full-preservation basin has mean defect $0.17$; the full-destruction basin has mean defect $0.64$.

\begin{center}
\begin{tabular}{lcccc}
\toprule
\textbf{Basin} & \textbf{Fraction} & \textbf{Mean defect} & \textbf{Final surviving} & \textbf{Mean $|E|$} \\
\midrule
Full preservation & 42\% & 0.17 & 14/14 & 14.6 \\
Partial & 36\% & 0.38 & 2--6/14 & 9.1 \\
Full destruction & 22\% & 0.64 & 0/14 & 5.7 \\
\bottomrule
\end{tabular}
\end{center}

\textbf{Finding 3: Structural feedback loop.} Trajectory tracing reveals a self-reinforcing mechanism. In the preservation basin, the fraction of defect-free extensions \emph{increases} along the trajectory (from $22/28 = 79\%$ to $33/36 = 92\%$): graphs that gain edges (via radius-2 local additions) develop richer local structure that stabilizes refinements. In the destruction basin, a single defect-producing step drops the defect-free fraction to near zero ($1/18$), trapping the trajectory. The first few steps determine the basin: early defect~$\leq 0.0$ predicts $66\%$ final survival; early defect~$> 0.5$ predicts $29\%$ final survival.

\textbf{Finding 4: Sharp Boltzmann phase transition.} Even mild selective pressure ($\beta = 0.25$, corresponding to a ${\sim}22\%$ multiplicative preference for defect-free extensions) collapses mean defect from $0.37$ to $0.11$ and raises full-preservation probability from $45\%$ to $81\%$:

\begin{center}
\begin{tabular}{cccc}
\toprule
$\boldsymbol{\beta}$ & \textbf{Mean defect} & \textbf{$P(\text{full pres.})$} & \textbf{$P(\text{defect}{<}0.3)$} \\
\midrule
0.00 & 0.37 & 45\% & 50\% \\
0.25 & 0.11 & 81\% & 82\% \\
0.50 & 0.03 & 91\% & 95\% \\
1.00 & 0.03 & 89\% & 99\% \\
$\geq 2$ & $< 0.03$ & $\sim 90\%$ & 100\% \\
\bottomrule
\end{tabular}
\end{center}

The transition is first-order-like: a discontinuous jump in preservation probability between $\beta = 0$ and $\beta = 0.25$, with saturation by $\beta \approx 0.5$. This means the DSP does not require strong enforcement---weak structural preference suffices.

\textbf{Kernel-level statement:} The Defect Selection Principle is \emph{partially emergent} on the barbell substrate. Three structural mechanisms support it without independent postulation: (a)~landscape pre-adaptation ($79\%$ of local extensions are defect-free); (b)~bimodal attractor basins ($42\%$ of random trajectories fully preserve all refinements); (c)~sharp Boltzmann phase transition (mild selective pressure $\beta \geq 0.25$ is sufficient). Under purely uniform random dynamics, mean defect is $0.35$, not zero---the DSP is not automatic. But the configuration space geometry is \emph{pre-adapted}: the structural prerequisites are in place, and only weak enforcement is required to make the DSP effectively exact. This downgrades DSP's status from ``independent assumption'' to ``emergent with mild structural support.''


\subsection{Toy Model 7: Certified Cycle Supply and Dimensional Viability}\label{app:toymodel7}

\statusnotice{Computational verification of the witness-inequivalent variables and certified supply inequality D11$''$ across 25 graph families ($n = 4$--$16$), with three gadget families ($C_3$, $C_4$, $Q_3$), strict certification, and carrier-summed supply.}

\textbf{Purpose.} Raw gadget counts used in dimensional selection can be inflated by redundant copies and blinded by basis inadequacy. This toy model implements the full provenance-discharged canonical variables---witness-inequivalent overlap $M^\perp$, carrier-summed supply $\tilde\sigma^\perp$, and strict certification fractions $\phi_i$---and tests the certified supply inequality on explicit substrates.

\textbf{Definitions (kernel-licensed).}

\emph{Gadget families:} $G_1 = C_3$ (induced triangles), $G_2 = C_4$ (induced squares), $G_3 = Q_3$ (induced cubes).

\emph{Witness equivalence.} Two gadget instances sharing a carrier are witness-equivalent ($g \sim_W h$) if no admissible local test within radius $R_W$ can distinguish them. The witness-inequivalent multiplicity at carrier $c$ is $M_i^\perp(c) := |\mathcal{G}_i(c)/{\sim_W}|$.

\emph{Carrier-summed supply.} $\Sigma_i^\perp(U) := \sum_{c \in \mathrm{Car}_i(U)} M_i^\perp(c)$, normalized by the fixed number of carriers per gadget: $\tilde\sigma_1^\perp = \Sigma_1^\perp / 3$, $\tilde\sigma_2^\perp = \Sigma_2^\perp / 4$, $\tilde\sigma_3^\perp = \Sigma_3^\perp / 6$. No tunable constants are introduced.

\emph{Strict certification.} A triangle certifies a fundamental cycle iff its 3 edges are all contained in the cycle. A square certifies iff its 4 edges are contained. A cube certifies iff at least one face-square has all 4 edges contained. Coverage fractions: $\phi_i(U) :=$ fraction of BFS fundamental cycles strictly certified by gadget type $i$.

\textbf{Certified Supply Inequality (D11$''$).} The witness-canonical supply bound:
\[
\beta(U) \;\leq\; \frac{\tilde\sigma_1^\perp(U)}{\phi_1(U)} + \frac{\tilde\sigma_2^\perp(U)}{\phi_2(U)} + \frac{\tilde\sigma_3^\perp(U)}{\phi_3(U)} + o(|U|),
\]
with terms omitted when $\phi_i = 0$. This replaces all earlier raw-count requirement inequalities. It simultaneously addresses \emph{basis blindness} (if a regime is triangle-dominated, $\phi_2 = \phi_3 = 0$ and only the triangle term is available) and \emph{witness collapse} (carrier-summed supply $\tilde\sigma^\perp$ remains correct even when class counts $\sigma^\perp$ collapse under symmetry).

\textbf{Results.} We computed all variables on 25 graphs across 7 families. Selected results:

\begin{center}
\begin{tabular}{lccccccccc}
\toprule
\textbf{Family} & $n$ & $\beta$ & $\triangle$ & $\square$ & $\tilde\sigma_2^\perp$ & $\phi_1$ & $\phi_2$ & \textbf{D11$''$} & \textbf{Regime} \\
\midrule
Grid $r{\times}c$ & 6--16 & 2--9 & 0 & 2--9 & 1.75--9.0 & 0 & 0.33--1.0 & 19/21$^*$ & 2-cell \\
Ladder-$\ell$ & 6--12 & 2--5 & 0 & 2--5 & 1.75--4.5 & 0 & 0.60--1.0 & 2/4$^*$ & 2-cell \\
$Q_3$ (3-cube) & 8 & 5 & 0 & 6 & 3.0 & 0 & 0.80 & $\checkmark$ & 3-cell \\
Cyl-$2{\times}4$ & 8 & 5 & 0 & 6 & 3.0 & 0 & 0.80 & $\checkmark$ & 3-cell \\
Tor-$4{\times}4$ & 16 & 17 & 0 & 24 & 8.0 & 0 & 0.65 & $\checkmark$ & 3-cell \\
\midrule
$K_3$-$b_N$-$K_3$ & 7--11 & 2 & 2 & 0 & 0 & 1.00 & 0 & $\checkmark$ & tri \\
$K_4$-$b_N$-$K_4$ & 9--11 & 6 & 8 & 0 & 0 & 1.00 & 0 & $\checkmark$ & tri \\
$K_4$-lad-$K_4$ & 12--14 & 9--10 & 8 & 3--4 & 3.0--3.75 & 0.40--0.44 & 0 & $\checkmark$ & tri \\
\bottomrule
\end{tabular}
\end{center}
{\footnotesize $^*$Small finite-size failures where $\tilde\sigma_2^\perp$ falls just below $\beta$; absorbed by the $o(|U|)$ slack in D11$''$.}

\textbf{Finding 1: Carrier-summed supply resolves witness collapse.} On the $4{\times}4$ torus, 24 raw squares collapse to $\sigma_2^\perp = 1$ witness class---but $\tilde\sigma_2^\perp = 8.0$ (the carrier-summed supply). This is the correct quantity for cycle-supply accounting: the torus has 8 units of certified square supply despite having only 1 witness class. The earlier $\sigma^\perp$-based D9 falsely rejected the torus; D11$''$ correctly accepts it.

\textbf{Finding 2: Triangles rescue all barbells.} Under the square/cube-only inequality, all barbells fail ($\phi_2 = \phi_3 = 0$). Adding triangles ($G_1 = C_3$) with carrier-summed supply gives $\tilde\sigma_1^\perp / \phi_1 \geq \beta$ for \emph{all} barbells tested (including $K_4$-barbells). The carrier-summed supply is essential: $K_4$-barbells have $\sigma_1^\perp = 3$ (only 3 witness classes) but $\tilde\sigma_1^\perp = 6.0$ (each clique edge carries $M_1^\perp \geq 1$). Eight graphs are rescued by triangles (D9 fails, D11$''$ passes).

\textbf{Finding 3: D11$''$ passes on 21/25 graphs.} The 4 remaining failures are all small finite-size effects (Grid-$2{\times}3$, Ladder-3, Ladder-5, Cyl-$2{\times}6$) where $\tilde\sigma_2^\perp$ falls within 0.25--0.70 of $\beta$---well within the $o(|U|)$ boundary discount. On all sufficiently large graphs ($n \geq 8$, connected, $\beta \geq 3$), D11$''$ is satisfied.

\textbf{Finding 4: Basis adequacy trichotomy.} Strict certification fractions classify substrates into three structural universality classes:
\begin{itemize}[nosep]
  \item \textbf{Triangle-dominated} ($\phi_1 > 0$, $\phi_2 = \phi_3 = 0$): all barbells, $K_4$-ladder-$K_4$. Only triangle-supported response invariants are admissible.
  \item \textbf{2-cell} ($\phi_2 > 0$, $\phi_3 = 0$): grids, ladders, Cyl-$2{\times}6$. Square-supported invariants available.
  \item \textbf{3-cell} ($\phi_3 > 0$): $Q_3$, Cyl-$2{\times}4$, Cyl-$3{\times}4$, Tor-$3{\times}4$, Tor-$4{\times}4$. Cube-supported invariants available; response algebra dimension strictly larger.
\end{itemize}
This classification is not geometric---it is a certification classification of cycle-space visibility.

\textbf{Finding 5: $\tau$--D11$''$ tension is resolved.} In the earlier D9 formulation, $\tau$-optimal graphs (barbells) failed the requirement inequality and D9-optimal graphs (grids) had low $\tau$. Under D11$''$, \emph{all barbells pass} via the triangle term. The structural tension survives only as a question about which certification class a given substrate falls into---not whether it can support dimensional selection at all. The $K_4$-ladder-$K_4$ family occupies a transitional position: it passes D11$''$ via triangles while also containing squares.

\textbf{Dimensional Viability Band (D12$''$).} Combining D11$''$ with overlap ceilings $\Lambda_{i,\max}$ yields feasibility bounds $\tilde\sigma_i^\perp(U) \leq \mathrm{Cap}_i(\Lambda_{i,\max}) \cdot |U|$. In any regime where $\beta(U) = \Theta(|U|)$, viability requires
\[
1 \;\lesssim\; \frac{\mathrm{Cap}_1(\Lambda_{1,\max})}{\phi_1} + \frac{\mathrm{Cap}_2(\Lambda_{2,\max})}{\phi_2} + \frac{\mathrm{Cap}_3(\Lambda_{3,\max})}{\phi_3}.
\]
Overlap ceilings (defect pressure) constrain capacity, which constrains which certification class can sustain $\beta = \Theta(|U|)$. This is the dimensional selection mechanism stated without geometry.

\textbf{Certification-gated response laws.} A gadget family with $\phi_i = 0$ cannot generate stable response invariants: its cycles are invisible to the certified supply, so no compression-stable boundary functional can depend on it. This gives the precise analogue of Lovelock's dimensional gating: only certification-visible gadget families produce admissible response-law terms. In the 2-cell class, cube terms vanish; in the triangle class, square and cube terms vanish. The response algebra dimension increases monotonically with the visible gadget set, and strictly so under \emph{non-redundant visibility} (there exist collar types with equal square supply but different cube supply) and \emph{spectral activation} (the cube feature direction has nonzero projection onto the stable eigenspace of the coarse-graining operator).

\textbf{Kernel-level summary.} The canonical variables $M^\perp$, $\tilde\sigma^\perp$, $\phi$ are computationally well-defined, and the certified supply inequality D11$''$ passes on 21/25 tested graphs (all failures are small finite-size effects). Carrier-summed supply resolves witness collapse; triangles resolve basis blindness. The dimensional selection mechanism is now: cycle demand $\leq$ certified supply $\leq$ overlap-constrained capacity. No geometric primitives are assumed at any stage.


\subsection{Toy Model 8: Ladder Certificates and the Dimensional Response Jump}\label{app:toymodel8}

\statusnotice{Computational implementation of the Ladder Certificate (LC) protocol. First certified strict dimensional ladder jumps: the response-law algebra dimension increases when cube structure becomes visible.}

\textbf{Purpose.} Toy Model~7 establishes the certified supply inequality D11$''$ and the basis adequacy trichotomy. The deeper question is whether transitioning from a 2-cell regime ($\phi_3 = 0$) to a 3-cell regime ($\phi_3 > 0$) forces a \emph{strict increase} in the dimension of the admissible response-law algebra---the NFC analogue of ``a new Lovelock invariant becomes available in higher dimension.'' The Ladder Certificate protocol makes this question mechanical and referee-auditable.

\textbf{Response-law algebra.} For a region family $U_n$, define collar types $\mathcal{T}^\perp$ as rooted witness-canonical descriptors at boundary sites (multisets of $(d, \deg, \Delta_i, s_i)$ tuples within a $k^*$-ball). The type-level random walk matrix $T^\perp$ captures how collar-type frequencies transform under WL-partition coarse-graining $\Pi_b$. The compression-stable response algebra $\mathscr{E}^{\text{obs}}_{\boldsymbol\phi}$ consists of boundary functionals that are left eigenvectors of $T^\perp|_S$, restricted to the support set $S$ of types with positive limiting frequency. Its dimension counts the number of independent response laws that survive coarse-graining.

\textbf{Strict ladder theorem (D.4/D.6).} The observable response algebra dimension satisfies $\dim(\mathscr{A}^{\text{obs}}_{\boldsymbol\phi}(\mathcal{R}_3)) > \dim(\mathscr{A}^{\text{obs}}_{\boldsymbol\phi}(\mathcal{R}_2))$ whenever three conditions hold:
\begin{enumerate}[nosep]
  \item \textbf{Cube visibility:} $\phi_3 > 0$ in $\mathcal{R}_3$.
  \item \textbf{Non-redundant visibility (NonRed):} there exist collar types $t, t' \in S_3$ with $g_2(t) = g_2(t')$ but $g_3(t) \neq g_3(t')$, where $g_i$ is the gadget-$i$ incidence feature map.
  \item \textbf{Spectral activation (SpecAct):} the projection of the $g_3$ direction onto the stable eigenspace of $T^\perp|_{S_3}$ has nonzero norm: $\eta = \|\mathrm{Proj}_{E^{(3)}}(g_3)\| > 0$.
\end{enumerate}
NonRed ensures the cube feature creates genuinely new observables (not reducible to square features); SpecAct ensures these survive coarse-graining.

\textbf{Lemma D.7 (NonRed is necessary for cube-driven strict growth).}
Let $\mathscr{A}_2 := \mathrm{Alg}\langle 1, g_2\rangle$ and $\mathscr{A}_{2,3} := \mathrm{Alg}\langle 1, g_2, g_3\rangle \subseteq \mathbb{R}^{S_3}$.
If NonRed fails on~$S_3$ (i.e., $g_2(t)=g_2(t') \Rightarrow g_3(t)=g_3(t')$), then $\mathscr{A}_{2,3} = \mathscr{A}_2$.
\emph{Proof.}
NonRed failure means $g_3$ is constant on each fiber of~$g_2$, so there exists $h\colon g_2(S_3)\to\mathbb{R}$ with $g_3(t)=h(g_2(t))$ for all~$t$.
Since $g_2(S_3)$ is finite, Lagrange interpolation yields a polynomial $p$ with $h(x)=p(x)$ on~$g_2(S_3)$.
Hence $g_3 = p \circ g_2 \in \mathscr{A}_2$, so $\mathscr{A}_{2,3} = \mathscr{A}_2$.\quad$\square$

\textbf{Corollary D.8 (Necessity for compression-stable strict ladder under cube-only novelty).}
If $S_2\subseteq S_3$, the observable algebra on~$S_3$ is generated by $\{1,g_2,g_3\}$ with no additional independent collar-type features beyond those on~$S_2$, and NonRed fails, then $\dim(\mathscr{E}^{\text{obs}}_{\boldsymbol\phi}(\mathcal{R}_3)) = \dim(\mathscr{E}^{\text{obs}}_{\boldsymbol\phi}(\mathcal{R}_2))$.
Thus any strict ladder jump attributable to cube structure requires NonRed.

\textbf{Lemma D.9 (SpecAct is automatic under primitivity).}
If the restricted operator $T_3 := T^\perp|_{S_3}$ is primitive (irreducible and aperiodic) and some $t\in S_3$ has $g_3(t) > 0$, then SpecAct holds: $\eta > 0$.
\emph{Proof.}
By Perron--Frobenius, $T_3$ has a unique stationary distribution~$\pi$ with $\pi(t) > 0$ for all $t\in S_3$, and $\pi^\top \in E^{(3)}$.
Then $\langle \pi, g_3\rangle = \sum_t \pi(t)\,g_3(t) > 0$ since $\pi > 0$ and $g_3 \geq 0$ with $g_3(t) > 0$ for some~$t$.
Hence $\mathrm{Proj}_{E^{(3)}}(g_3) \neq 0$.\quad$\square$

\textbf{Corollary D.10 (SpecAct reduces to support-level cube visibility).}
Under the hypotheses of Lemma~D.9, SpecAct is equivalent to the existence of a supported collar type with $g_3 > 0$.
In primitive regimes, SpecAct adds no independent constraint beyond $\phi_3 > 0$ together with the definition of~$S_3$ as the positive-frequency support.

\emph{Synthesis (strict ladder gates, minimized).}
Under primitivity of $T^\perp|_{S_3}$, the strict ladder reduces to
\[
\phi_3 > 0 \;\text{and}\; \mathrm{NonRed} \quad\Longrightarrow\quad
\dim\!\big(\mathscr{E}^{\text{obs}}(\mathcal{R}_3)\big) > \dim\!\big(\mathscr{E}^{\text{obs}}(\mathcal{R}_2)\big),
\]
with SpecAct subsumed by Lemma~D.9. By Lemma~D.7/Corollary~D.8, NonRed is also necessary. Hence: the strict ladder holds if and only if $\phi_3 > 0$ and NonRed, given primitivity.

\textbf{Results.} We tested the LC protocol on 12 regime pairs. Two substrates achieve full strict ladder certification:

\begin{center}
\begin{tabular}{lccccccc}
\toprule
\textbf{R3 substrate} & $\phi_3$ & $|S_3|$ & \textbf{NonRed} & \textbf{Prim} & \textbf{SpecAct} & $d_2$ & $d_3$ \\
\midrule
Fan3+$Q_3$ & 0.50 & 7 & $\checkmark$ & $\checkmark$ & $\checkmark$ ($\eta = 1.73$) & 3 & 7 \\
Grid$_{3\times3}$+Fan3+$Q_3$ & 0.33 & 12 & $\checkmark$ & $\checkmark$ & $\checkmark$ ($\eta = 1.73$) & 3 & 11 \\
\midrule
\multicolumn{8}{l}{\small\emph{NonRed-failing substrates (for comparison):}} \\
Grid$_{4\times4}$+$Q_3$(bridge) & 0.29 & 9 & $\times$ & $\checkmark$ & $\checkmark$ & 3 & 9 \\
Cyl-$3{\times}4$ & 0.44 & 2 & $\times$ & $\checkmark$ & $\checkmark$ & 3 & 2 \\
$Q_3$ alone & 0.80 & 1 & $\times$ & $\checkmark$ & $\checkmark$ & 3 & 1 \\
\bottomrule
\end{tabular}
\end{center}

The R2 baseline is Grid-$4{\times}4$ ($d_2 = 3$, $\phi_2 = 0.33$, $\phi_3 = 0$) in all cases.

\smallskip
\textbf{Explicit Ladder Certificate: LC(Grid-$4{\times}4$ $\to$ Fan3+$Q_3$).}
The collar-type alphabet $\mathcal{T}^\perp$ at $k^*\!=\!1$ uses rooted multisets of
$(\mathrm{dist},\deg,\deg_{\mathrm{ball}},\triangle\text{-inc},\square\text{-inc},Q_3\text{-inc})$ tuples
within the radius-1 ball.  Fan3+$Q_3$ has $|S_3| = 7$ collar types; we display the
root-vertex features and the feature maps $g_2,g_3$:

\begin{center}\small
\begin{tabular}{clcc}
\toprule
idx & root-vertex collar label & $g_2$ & $g_3$ \\
\midrule
0 & $[\deg\!=\!2,\;\triangle\!=\!0,\;\square\!=\!0,\;Q_3\!=\!0,\;|\text{ball}|\!=\!3]$ & 0.00 & 0.00 \\
1 & $[\deg\!=\!2,\;\triangle\!=\!0,\;\square\!=\!1,\;Q_3\!=\!0,\;|\text{ball}|\!=\!3]$ & 1.00 & 0.00 \\
2 & $[\deg\!=\!2,\;\triangle\!=\!0,\;\square\!=\!1,\;Q_3\!=\!0,\;|\text{ball}|\!=\!3]$ & 1.00 & 0.00 \\
3 & $[\deg\!=\!7,\;\triangle\!=\!0,\;\square\!=\!3,\;Q_3\!=\!0,\;|\text{ball}|\!=\!8]$ & 3.00 & 0.00 \\
4 & $[\deg\!=\!3,\;\triangle\!=\!0,\;\square\!=\!3,\;Q_3\!=\!1,\;|\text{ball}|\!=\!4]$ & 3.00 & 1.00 \\
5 & $[\deg\!=\!3,\;\triangle\!=\!0,\;\square\!=\!3,\;Q_3\!=\!1,\;|\text{ball}|\!=\!4]$ & 3.00 & 1.00 \\
6 & $[\deg\!=\!4,\;\triangle\!=\!0,\;\square\!=\!3,\;Q_3\!=\!1,\;|\text{ball}|\!=\!5]$ & 3.00 & 1.00 \\
\bottomrule
\end{tabular}
\end{center}

\noindent\textbf{NonRed witness pair.}  Types $t = \text{idx}\;4$ and $t' = \text{idx}\;3$ satisfy
$g_2(t) = g_2(t') = 3.00$ but $g_3(t) = 1.00 \neq 0.00 = g_3(t')$.
Their full collar signatures are:

\smallskip\noindent
$t$: $\;[\deg\!=\!3,\;\square\!=\!3,\;Q_3\!=\!1,\;|\text{ball}|\!=\!4]$
--- a degree-3 cube vertex whose $k^*$-ball contains 4 vertices, all with
$(\square\text{-inc},Q_3\text{-inc}) = (3,1)$.  This is a generic interior vertex of $Q_3$.

\smallskip\noindent
$t'$: $\;[\deg\!=\!7,\;\square\!=\!3,\;Q_3\!=\!0,\;|\text{ball}|\!=\!8]$
--- the fan center, degree~7 (3 fan arms $+$ bridge $+$ 3 return edges),
whose $k^*$-ball contains 8 vertices, all with $Q_3$-incidence zero.
It participates in exactly 3 induced squares (one per fan arm) and no cubes.

\smallskip\noindent
\emph{Structural interpretation.}  Both $t$ and $t'$ participate in exactly 3 squares,
so the square feature $g_2$ cannot distinguish them.  But $t$ lies inside a cube
while $t'$ does not, so $g_3$ separates them.  This is the defining signature of
non-redundant cube visibility: the cube feature carries information not reducible to
the square feature on the observed support.

\smallskip\noindent
\textbf{SpecAct:} $\eta = \|\mathrm{Proj}_{E^{(3)}}(g_3)\| = 1.73$. \textbf{PASS.}

\noindent
\textbf{DIM:} $T_2$ spectrum $\{1.00, 0.33, -0.50\}$, $d_2 = 3$.\\
$T_3$ spectrum $\{1.00, 0.97, 0.42, 0.06, -0.24, -0.72, \ldots\}$, $d_3 = 7$.\\
Strict jump $d_3 = 7 > 3 = d_2$.  \textbf{All gates pass; strict ladder certified.}

\textbf{Finding 1: First certified strict ladder jumps.} Fan3+$Q_3$ achieves $d_2 = 3 \to d_3 = 7$ (response algebra dimension more than doubles). Grid$_{3\times3}$+Fan3+$Q_3$ achieves $d_2 = 3 \to d_3 = 11$. Both pass all three gates. The jump is genuine: the 3-cell regime supports strictly more compression-stable response invariants than the 2-cell regime.

\textbf{Finding 2: NonRed is the binding constraint.} Of 12 regime pairs tested, all with $\phi_3 > 0$ pass SpecAct, but only 2 pass NonRed. The explicit witness pair above shows the mechanism: the fan center ($t' = \text{idx}\;3$, $\deg = 7$, $\square = 3$, $Q_3 = 0$) and a cube vertex ($t = \text{idx}\;4$, $\deg = 3$, $\square = 3$, $Q_3 = 1$) share identical square incidence $g_2 = 3$ but differ on cube incidence $g_3$. On pure grids bridged to $Q_3$, every vertex with $g_2 = 3$ also has $g_3 = 1$, making $g_3$ functionally dependent on $g_2$ and NonRed impossible. The ``square fan'' construction breaks this correlation by placing a non-cube vertex at exactly the right square count.

\textbf{Finding 3: NonRed is the NFC analogue of functional independence.} In Lovelock's framework, dimension $d$ produces a new invariant (the Gauss--Bonnet term in $d = 4$) precisely because the Riemann curvature carries information not reducible to the Ricci tensor. NonRed is the exact NFC translation: the cube feature $g_3$ must carry information not reducible to the square feature $g_2$ on the observed collar-type support. Substrates where cube vertices are ``too similar'' to square vertices fail this test---the cube structure exists but adds no new observable.

\textbf{Finding 4: SpecAct is generically satisfied (explained).} In all tested 3-cell regimes, $T^\perp|_{S_3}$ is primitive (Prim column, all $\checkmark$) and $S_3$ contains cube-bearing collar types ($g_3 > 0$), hence SpecAct holds automatically by Lemma~D.9. SpecAct failure would require either non-primitivity of $T^\perp|_{S_3}$ (reducible or periodic type dynamics) or a mismatch between $\phi_3 > 0$ and support-level cube incidence. Neither is observed on any tested substrate. By Corollary~D.10, under primitivity SpecAct reduces to $\phi_3 > 0$, so the effective strict ladder condition is simply $\phi_3 > 0$ and NonRed.

\textbf{Finding 5: Symmetry suppresses NonRed (Proposition D.11).}
The rarity of NonRed has a structural explanation: graph automorphisms force functional dependence between $g_2$ and $g_3$.

\emph{Proposition D.11 (Symmetry constrains NonRed via fiber transitivity).}
Let $O_1,\ldots,O_k$ be the orbits of $\mathrm{Aut}(G)$ on $V(G)$.
(i)~Both $g_2$ and $g_3$ are $\mathrm{Aut}(G)$-invariant, hence constant on each orbit.
(ii)~If $G$ is vertex-transitive ($k=1$), NonRed is \emph{impossible} (one orbit $\Rightarrow$ every $g_2$-fiber has constant $g_3$).
(iii)~In general, NonRed requires a \emph{split $g_2$-fiber}: orbits $O_i, O_j$ with $g_2(O_i) = g_2(O_j)$ but $g_3(O_i) \neq g_3(O_j)$.
The critical point is not $|\mathrm{Aut}(G)|$ per se, but whether the orbit partition separates $g_3$ values within some $g_2$-level set. The mechanism is confirmed by exact automorphism counts: Fan3+$Q_3$ has $|\mathrm{Aut}| = 288$ yet passes NonRed (the fan center and cube vertices lie in different orbits at $g_2 = 3$), while Fan4+$Q_3$ has $|\mathrm{Aut}| = 2304$ and fails (fan center has $g_2 = 4 \neq 3$, so no fiber overlaps).

The mechanism is visible in the $g_2$-fiber structure:

\begin{center}\small
\begin{tabular}{lcl}
\toprule
\textbf{Substrate} & $g_2 = 3$ \textbf{fiber} & \textbf{NonRed} \\
\midrule
$Q_3$ (vertex-trans., $|\mathrm{Aut}| = 48$) & $\{g_3 = 1\}$ & fail \\
Cyl-$2{\times}4$ (vertex-trans., $|\mathrm{Aut}| = 48$) & $\{g_3 = 1\}$ & fail \\
Grid$_{3\times3}$+$Q_3$(b1) & $\{g_3 = 1\}$ & fail \\
Grid$_{4\times4}$+$Q_3$(b1) & $\{g_3 = 1\}$ & fail \\
Fan4+$Q_3$ & $\{g_3 = 1\}$ & fail \\
Fan3+$Q_3$ & $\{g_3 = 0,\; g_3 = 1\}$ & \textbf{PASS} \\
Grid$_{3\times3}$+Fan3+$Q_3$ & $\{g_3 = 0,\; g_3 = 1\}$ & \textbf{PASS} \\
\bottomrule
\end{tabular}
\end{center}

On all substrates except those containing the Fan3 construction, the $g_2 = 3$ fiber is a singleton $\{g_3 = 1\}$: every vertex with 3 square memberships is also in a cube. The fan center (degree 7, $\square = 3$, $Q_3 = 0$) is the \emph{only} tested structure that places a non-cube vertex at $g_2 = 3$, splitting the fiber. This is an arithmetic coincidence: each fan arm contributes one induced $C_4$, and $Q_3$ vertices each lie in exactly 3 induced $C_4$s, so only $n_{\mathrm{fan}} = 3$ achieves the match.

\emph{Structural conclusion.} The dimensional ladder is a \emph{rarity condition}. Symmetric substrates generically force $g_3 = f(g_2)$ on support, preventing NonRed. This is the NFC analogue of the observation that symmetric spacetimes (de~Sitter, FLRW) have too few independent curvature components for higher Lovelock invariants to be nontrivial; generic spacetimes break this constraint.

\textbf{Interpretation.} The dimensional ladder is not automatic ($\phi_3 > 0$ is necessary but not sufficient), not impossible (certified on concrete substrates), and structurally rare (suppressed by symmetry). By Lemma~D.7, NonRed is \emph{necessary} for cube-driven algebra enlargement; by D.4/D.6 it is (with $\phi_3 > 0$) \emph{sufficient} under primitivity (Lemma~D.9). By Proposition~D.11, graph symmetry generically forces $g_3$ to be functionally dependent on~$g_2$, killing NonRed. The sharp characterization is:
\begin{multline*}
\text{strict dimensional ladder} \;\iff\; \phi_3 > 0 \;\text{and}\; \mathrm{NonRed} \\
\text{(under primitivity of } T^\perp|_{S_3}\text{)}.
\end{multline*}
The complete provenance chain is now:
\begin{multline*}
\text{cycle demand} \;\leq\; \text{certified supply} \;\leq\; \text{overlap-constrained capacity} \\
\;\Rightarrow\; \text{certification-gated response algebra} \;\Rightarrow\; \text{strict ladder under NonRed + SpecAct}.
\end{multline*}

No step assumes geometric embedding or manifold structure. The dimensional selection mechanism, the Lovelock analogue, and the dimensional ladder are all structurally locked within the NFC kernel plus the explicitly declared coarse-graining scheme (WL-1 partition, which is canonical under isomorphism and uses no coordinates).


\subsection{Synthesis and Implications}

Across the eight computational models (3--35 vertices; unary, binary, phase-valued, systematic multi-class, dynamical, witness-inequivalent, and ladder-certificate signatures):
\begin{itemize}[nosep]
  \item Behavioral congruence equals isomorphism (Theorem~\ref{thm:behav-iso}); the admissible symmetry group $G$ is non-transitive in general (Toy Model~1), with cross-orbit weighting as empirical input (Remark~\ref{rem:transitivity}).
  \item Projectable refinements are strongly selected toward global uniformity ($r_1$/$r_6$ survive indefinitely); local, unbalanced, directed, periodic, or multi-valued distinctions are erased under natural extensions.
  \item The refinement ceiling $\tau$ depends critically on the strength of the admissible class: $\tau = 0$ under global operations (complement, delete-all, complete-all), but grows without bound under locality-respecting extensions (Toy Model~5). Bounded propagation is a structural prerequisite.
  \item The $\tau = 4$ threshold required by Theorem~\ref{thm:ceiling} is first achieved at $n = 9$ on the connected barbell $K_3\text{-}b_3\text{-}K_3$. The barbell family provides a concrete substrate for the dimensional selection mechanism with $\tau$ scaling linearly with bridge length.
  \item Under multi-arity refinements (vertex, edge, path, triangle, and 4-cycle orbit types), the local load reaches $g_{\min} = 30$ on $K_4$-ladder-$K_4$ structures and $g_{\max} > 32$ on clique-adjacent vertices. The remaining gap at bridge midpoints requires minimum degree $\geq 3$, a natural structural condition.
  \item The Defect Selection Principle is partially emergent (Toy Model~6): the configuration space geometry of $K_3\text{-}b_3\text{-}K_3$ is pre-adapted for low-defect dynamics (79\% of local extensions are defect-free), with bimodal attractor basins and a sharp Boltzmann phase transition at $\beta \approx 0.25$. DSP's status is downgraded from ``independent assumption'' to ``emergent with mild structural support.''
  \item The canonical witness-inequivalent variables $M^\perp$, $\tilde\sigma^\perp$, $\phi$ (Toy Model~7) reveal massive redundancy in raw gadget counting (up to $24\times$ on tori) and, after three audit-driven corrections (strict certification, carrier-summed supply, triangle inclusion), yield the certified supply inequality D11$''$: $\beta \leq \sum_i \tilde\sigma_i^\perp / \phi_i + o(|U|)$. D11$''$ passes on 21/25 tested graphs; all 4 failures are small finite-size effects. Adding triangles as a gadget family rescues all barbells (D9 fails, D11$''$ passes via $\phi_1 > 0$). The basis adequacy trichotomy (tri-dominated / 2-cell / 3-cell) classifies substrates by cycle-space visibility, not geometry. The earlier $\tau$--D9 tension is resolved: all $\tau$-optimal barbells now satisfy D11$''$.
  \item The Ladder Certificate protocol (Toy Model~8) produces the first certified strict dimensional ladder jumps: $d_2 = 3 \to d_3 = 7$ (Fan3+$Q_3$) and $d_2 = 3 \to d_3 = 11$ (Grid$_{3\times3}$+Fan3+$Q_3$), with explicit collar-type witness pairs enumerated. NonRed is both necessary (Lemma~D.7: without it, $g_3 \in \mathrm{Alg}\langle g_2\rangle$) and, with $\phi_3 > 0$, sufficient (D.4/D.6). SpecAct is automatic under primitivity of $T^\perp|_{S_3}$ (Lemma~D.9, verified on all tested 3-cell substrates). The effective characterization is: strict ladder $\iff$ $\phi_3 > 0$ and NonRed (given primitivity). Graph symmetry generically suppresses NonRed by forcing $g_3$-constancy within $g_2$-fibers (Proposition~D.11); the ladder is a structural rarity condition, the NFC analogue of symmetric spacetimes having too few independent curvature components for higher Lovelock terms.
\end{itemize}

These results are unconditional within the kernel and rely solely on finite structures, explicit admissible classes, and the definitions of behavioral congruence and projectability. They provide concrete support for the unavoidability of compression (Theorem~\ref{thm:compress}), entropy monotonicity, finite refinement ceilings (Theorem~\ref{thm:finite}), direct evidence that the dimensional selection mechanism (Theorem~\ref{thm:ceiling}) can engage on connected structures of modest size, and---through the certified supply inequality D11$''$ and the Ladder Certificate protocol---a complete, provenance-audited chain from kernel axioms through dimensional selection to strict dimensional ladder jumps. No geometric embedding or manifold structure is assumed at any stage; the coarse-graining scheme (WL-1 partition) is the only declared combinatorial choice.

\begin{remark}[Computational Verification of Phase-A Hypotheses]\label{rem:preadapt-verify}
The Phase-A structural hypotheses (Assumptions~\ref{ass:preadapt},~\ref{ass:LS0}, Definition~\ref{def:LS2}) were computationally verified on all tested substrates:

\emph{Pre-adaptation.} For each substrate, local extension neighborhoods $\mathcal{N}(\tilde{\sigma})$ were enumerated (single-edge add/remove within radius~2), and the step defect $d(E;\tilde{\sigma})$ and low-defect fraction $\alpha_\varepsilon$ computed.
\begin{center}\small
\begin{tabular}{lrrcccc}
\toprule
Substrate & $n$ & $|\mathcal{N}|$ & $\alpha_{0.1}$ & $\alpha_{0.0}$ & $\bar{d}$ & LS-2 \\
\midrule
Grid-$3{\times}3$ & 9 & 26 & 1.000 & 1.000 & 0.000 & $\checkmark$ \\
Grid-$3{\times}4$ & 12 & 39 & 1.000 & 1.000 & 0.000 & $\checkmark$ \\
$Q_3$ & 8 & 24 & 1.000 & 1.000 & 0.000 & $\checkmark$ \\
Cyl-$2{\times}4$ & 8 & 24 & 1.000 & 1.000 & 0.000 & $\checkmark$ \\
Fan3+$Q_3$ & 19 & 66 & 0.985 & 0.939 & 0.004 & $\checkmark$ \\
\bottomrule
\end{tabular}
\end{center}
All substrates satisfy $\alpha_{0.1} \geq 0.985$, confirming preservation-basin structure with $\varepsilon = 0.1$, $\eta \leq 0.015$. The Fan3+$Q_3$ substrate---the only tested graph certifying a strict dimensional ladder---has the lowest $\alpha_{0.0}$ (0.939), indicating that the structurally richer topology required for NonRed introduces mild defect sensitivity but remains firmly in the pre-adapted regime.

\emph{LS-2 pairwise obstruction sufficiency.} For each substrate, all pairs of distinct WL-1 color classes were tested for differential response to single-edge probes. All substrates pass: every pair of distinct classes admits a local probe producing an order-dependent obstruction. Vertex-transitive graphs ($Q_3$, Cyl-$2{\times}4$) trivially satisfy LS-2 with zero pairs to separate. Fan3+$Q_3$ (8 WL classes, 28 pairs) is the most stringent test; all 28 pairs separate.
\end{remark}


\section{Smuggling Audit}\label{app:smuggle}

\subsection{Core Kernel Audit}

\begin{center}
\begin{tabular}{lcccp{4.5cm}}
\toprule
\textbf{Primitive} & \textbf{Assumed?} & \textbf{Derived?} & \textbf{Smuggled?} & \textbf{Notes} \\
\midrule
Real numbers $\mathbb{R}$ & No & No & No & Use $\mathbb{N}$ only \\
Hilbert spaces & No & No & No & V7 introduces representationally \\
Manifolds & No & Emergent & No & Conditional on Thm~\ref{thm:ceiling} \\
Probability axioms & No & Diagnostic & No & Intra-orbit uniformity forced; never load-bearing for exclusion \\
Time & No & Yes & No & Ledger coordinate \\
Dimensionality $\tau{=}4$ & No & Conditional & No & Phase-A stability selection \\
Lorentz signature & No & Yes & No & Conditional on isotropy + cone \\
Quantum interference & No & Yes & No & Conditional on phase model \\
Born rule & No & Yes & No & V7 representational layer \\
Holography & No & Yes & No & Conditional \\
Cross-orbit weighting & Yes & No & Honest & Empirical input (Remark~\ref{rem:transitivity}) \\
Law selection & No & Yes & No & Algebraic; weight-free (Theorem~\ref{thm:law}) \\
Bounded propagation & No & Derived & Upgraded & Phase-A basins exclude global moves (Thm~\ref{thm:locality}); derived from pre-adaptation, not assumed \\
Defect selection (DSP) & No & Derived & Upgraded & Defect budget derived from pre-adaptation + mild selection (Thm~\ref{thm:budget}); pre-adaptation verified computationally ($\alpha_{0.1} \geq 0.985$ on all substrates) \\
Phase-A pre-adaptation & Yes & No & Honest & Structural hypothesis (basin + trap asymmetry); verified computationally; explicit failure gates (F1--F4) \\
Local separability & Weakened & Conditional & Upgraded & Full separability replaced by LS-0 + LS-2 hierarchy; LS-2 verified on all tested substrates (28/28 pairs on Fan3+$Q_3$) \\
Witness-inequivalent variables & No & Yes & No & $M^\perp$, $\tilde\sigma^\perp$, $\phi$ derived from witness equivalence and carrier counting (Toy Model~7); no new primitives. Three audit corrections applied (strict cert., carrier-sum, triangles) \\
Certified supply inequality D11$''$ & No & Diagnostic & Resolved & Passes 21/25 graphs; all failures are finite-size. Triangles rescue barbells. Basis adequacy trichotomy classifies substrates by certification visibility \\
Certification-gated response laws & No & Structural & New & Gadget families with $\phi_i = 0$ cannot produce stable response invariants. Response algebra dimension increases with visible gadget set (Lovelock analogue) \\
Strict dimensional ladder jump & No & Computational & Certified & NonRed+SpecAct certified on Fan3+$Q_3$ ($d_3\!=\!7$) and Grid+Fan+$Q_3$ ($d_3\!=\!11$); witness pair enumerated (idx 3/4). NonRed necessary (D.7), sufficient with $\phi_3$ (D.4/D.6), rare under symmetry (D.11) \\
\bottomrule
\end{tabular}
\end{center}

\subsection{V7 Representational Extension Audit}

The V7 extension introduces three explicit axioms (NF, RBA, PNA). A detailed smuggling audit was conducted prior to ratification:
\begin{itemize}[nosep]
  \item \textbf{NF:} Pass. Obstruction claim; no observers or dynamics introduced.
  \item \textbf{RBA:} Pass with anti-ontology clause. Representational codomain only.
  \item \textbf{PNA:} Pass with tethering clauses. Contexts correspond to maximal compatible witness decompositions in $\mathcal{W}$; non-contextuality is invariance under admissible refactorizations.
  \item \textbf{Lemma V7-L1:} Pass. No observers, no time, PASS correctly enforced.
  \item \textbf{Theorem V7-T1:} Pass conditional on PNA tethering.
\end{itemize}


\section{Experimental Discriminators}\label{app:predict}

Potential empirical signatures, conditional on the framework's conditional realizations holding in the physical regime:

\begin{center}\small
\begin{tabular}{p{3.2cm}p{2.2cm}p{5.2cm}p{3.2cm}}
\toprule
\textbf{Signature} & \textbf{Required layer} & \textbf{Falsification condition} & \textbf{Approximate scale} \\
\midrule
Discrete spectra & Conditional (Part~III) & No survivor-count quantization window and no detectable Lorentz-violation onset near predicted saturation bands. & Regimes approaching $\tau$-saturation. \\
\addlinespace
Holographic noise & Conditional (Part~V) & Measured boundary/bulk fluctuation spectrum incompatible with finite survivor counting in exceptionally projectable systems. & Boundary-dominated high-coherence systems. \\
\addlinespace
Dimensional transitions & Conditional (Part~III) & No statistically significant Phase~C-style reorganization across controlled structural loading transitions. & Mesoscopic-to-cosmological crossover candidates. \\
\addlinespace
Structural uncertainty & Conditional (Part~IV) & Violation of the predicted complementary-refinement lower bound $\delta_A \cdot \delta_B \geq |\Sigma|/\tau$ in certified compatible setups. & Interference-sensitive finite-configuration experiments. \\
\bottomrule
\end{tabular}
\end{center}

\section{Framework Hierarchy and Version History}\label{app:hierarchy}

The NFC framework has been developed through a series of versioned layers:
\begin{enumerate}
  \item \textbf{Core Kernel (Parts~I--II):} Finite relational incidence, admissible extensions, compression, behavioral congruence = isomorphism, admissible law selection (algebraic, weight-free), non-functoriality, ledger coordinate, time emergence. Unconditional.
  \item \textbf{Conditional Realizations (Part~III):} Dimensional selection, Lorentzian signature, quantum interference, holographic bounds. Require stated structural conditions.
  \item \textbf{V7 Representational Extension (Section~\ref{sec:born}):} NF, RBA, PNA. Non-retroactive. Born-type rule as unique quadratic measure. Requires $\dim \geq 3$.
  \item \textbf{Research Program (Part~VI):} Einstein equations, Navier--Stokes, Standard Model gauge structure, generation bounds. Open problems.
\end{enumerate}

\medskip\noindent\textbf{Governing principle:} No result at a higher layer may modify, strengthen, or retroactively justify any result at a lower layer. The firewall between layers is strict. All claims are tagged: unconditional, conditional, representational~(V7), or open.


\section{Part~III Dependency Stack and Circularity Audit}\label{app:depstack}

This section records the dependency graph for Part~III and identifies the minimal assumption stack required for the $\tau$-selection mechanism. All arrows indicate strict logical dependence. No result is used before its prerequisites are defined.

\subsection*{Kernel Base (Unconditional)}
\begin{enumerate}[label=\textbf{K\arabic*.}]
  \item Finite relational structures and configuration spaces (Definitions~\ref{def:sig}--\ref{def:config}).
  \item Structural equivalence / isomorphism (Definition~\ref{def:equiv}).
  \item Admissible extension class axioms: closure, congruence preservation, definability closure, effective finiteness (Definition~\ref{def:admiss}).
  \item Behavioral congruence, survivor space, compression (Definitions~\ref{def:behav},~\ref{def:survivor}).
  \item Induced survivor maps and non-functoriality witness concept (Definitions~\ref{def:induced},~\ref{def:defect}; Theorem~\ref{thm:nonfunc}).
\end{enumerate}

\subsection*{Phase-A Package (One Sufficient Mechanism)}
Pre-adaptation and mild selection provide one mechanism by which Phase-A conditions (floor-attainment) may be realized. They are \emph{not} logically required for $\tau$-selection.
\begin{enumerate}[label=\textbf{PA\arabic*.}]
  \item \textbf{Structural Pre-Adaptation} (Assumption~\ref{ass:preadapt}): Preservation basin exists with high $\alpha_\varepsilon$ and trap asymmetry.
  \item \textbf{Mild Defect Aversion} (Assumption~\ref{ass:mild}): Selection not adversarial w.r.t.\ defect score.
  \item \textbf{No Probe Starvation} (Assumption~\ref{ass:nps}): Canonical probes not permanently suppressed.
  \item[$\Rightarrow$] \textbf{Emergent defect budget} (Theorem~\ref{thm:budget}): $\overline{\mathcal{D}} \leq \varepsilon + O(\eta)$. \emph{Derived from PA1+PA2.}
  \item[$\Rightarrow$] \textbf{Effective locality} (Theorem~\ref{thm:locality}): Global moves excluded from basin. \emph{Derived from PA1 + global sensitivity.}
  \item[$\Rightarrow$] \textbf{Floor-attainment}: Slack small. \emph{Derived from PA1+PA2+PA3.}
\end{enumerate}

\subsection*{Certified Floor and $\tau$-Selection (No Dynamics Required)}
\begin{enumerate}[label=\textbf{S\arabic*.}]
  \item \textbf{LS-0} (Assumption~\ref{ass:LS0}): Canonical local probe family exists.
  \item \textbf{LS-2} (Definition~\ref{def:LS2}): Pairwise obstruction sufficiency $\Rightarrow$ certified obstruction count $N_{\mathrm{cert}}(U)$.
  \item \textbf{Recurrence certificate} (Definition~\ref{def:reccert}): Coverage epoch $T_{\mathrm{cov}}(U)$.
  \item[$\Rightarrow$] \textbf{Certified defect floor} (Definition~\ref{def:floor}): $D_{\mathrm{floor}}(U) = c_\star N_{\mathrm{cert}}/T_{\mathrm{cov}}$. \emph{Purely structural.}
  \item \textbf{Local branching} (Definition~\ref{def:branch}): $L(U) \geq g(k)$.
  \item[$\Rightarrow$] \textbf{Linear obstruction growth} (Theorem~\ref{thm:obstruct}): $\geq c_0(L-1)$ events.
  \item[$\Rightarrow$] \textbf{Refinement ceiling} $\tau \leq 4$ (Theorem~\ref{thm:ceiling}): Obstruction pressure exceeds floor + slack for $k \geq 5$.
\end{enumerate}

\subsection*{Circularity Check}
The $\tau$-selection chain (S1--S7) depends only on LS-0, LS-2, recurrence certification, and local branching---all purely structural. It does \emph{not} depend on pre-adaptation, defect aversion, or any selection principle. The Phase-A mechanism package (PA1--PA3) provides one sufficient route to floor-attainment but is not logically required by the ceiling theorem. No result at level S depends on any result at level PA. The dependency graph is acyclic.


\section{Assumption and Claim Registry}\label{app:registry}

This registry centralizes the current assumption/load-bearing claim surface so audits can be run without traversing the full text.

\subsection*{Core assumptions and role}
\begin{center}\small
\begin{tabular}{p{2.7cm}p{3.4cm}p{2.2cm}p{3.8cm}}
\toprule
\textbf{Item} & \textbf{Statement class} & \textbf{Layer} & \textbf{Primary use} \\
\midrule
Definition~\ref{def:admiss} & Admissible extension axioms (closure/congruence/definability/finiteness) & Kernel & Enables behavioral congruence, entropy monotonicity, and law constraints. \\
\addlinespace
Definition~\ref{def:local-class} & Locality-respecting admissible classes (radius-bounded support) & Conditional interface & Canonical locality condition for Toy Model~5 and diversity scaling statements. \\
\addlinespace
Assumption~\ref{ass:LS0} & Canonical probe family exists & Conditional & Entry point for certified floor construction. \\
\addlinespace
Definition~\ref{def:LS2} & Pairwise obstruction sufficiency & Conditional & Supplies $N_{\mathrm{cert}}(U)$ in the ceiling chain. \\
\addlinespace
Definition~\ref{def:reccert} & Coverage-epoch recurrence certificate & Conditional & Supplies $T_{\mathrm{cov}}(U)$ for certified floor. \\
\addlinespace
Axioms~\ref{ax:nf},~\ref{ax:rba},~\ref{ax:pna} & Representational extension inputs & V7 & Projectivization and Born-type measure theorem. \\
\bottomrule
\end{tabular}
\end{center}

\subsection*{Flagship results and status}
\begin{center}\small
\begin{tabular}{p{3.5cm}p{2cm}p{2.2cm}p{4.2cm}}
\toprule
\textbf{Result} & \textbf{Tag} & \textbf{Status} & \textbf{Dependency pointer} \\
\midrule
Behavioral congruence $=$ isomorphism (Theorem~\ref{thm:behav-iso}) & Unconditional & Proved & Definition~\ref{def:admiss}, Definition~\ref{def:behav}. \\
\addlinespace
Law Selection (Theorem~\ref{thm:law}) & Unconditional$^\dagger$ & Proved & Definitions~\ref{def:admiss-sym},~\ref{def:law} + extension-separability premise. \\
\addlinespace
Refinement ceiling selection (Theorem~\ref{thm:ceiling}) & Conditional & Proved conditional & LS-0, LS-2, recurrence certificate, branching/obstruction chain. \\
\addlinespace
Born-Type probability (Theorem~\ref{thm:born}) & Representational (V7) & Proved conditional & NF+RBA+PNA + Gleason-class theorem for $\dim \ge 3$. \\
\addlinespace
Einstein/Navier--Stokes/SM derivations (Part~VI) & Open & Not derived & Research targets only; no completion claim. \\
\bottomrule
\end{tabular}
\end{center}

{\footnotesize $^\dagger$Unconditional relative to the kernel plus the explicit orbit extension-separability condition stated in Theorem~\ref{thm:law}.}

\section{Structural Hardening Catalog}\label{app:hardening}

This section catalogs the structural upgrades applied to the theoretical framework, with status and provenance.

\begin{center}\small
\begin{tabular}{p{4.5cm}p{2.5cm}p{6.5cm}}
\toprule
\textbf{Upgrade} & \textbf{Status} & \textbf{Effect} \\
\midrule
(iii) $\to$ (iii$'$): Definability Closure & \textbf{Integrated} & ``No new invariant partitions'' formalized as FO-definability over base signature. Prevents partition inflation; aligns with $\mathrm{Aut}(\mathfrak{A})$-invariance. \\
\addlinespace
Phase-A Pre-Adaptation & \textbf{Integrated} & Hard defect budget ($D_{\max}$) replaced by preservation basin + trap asymmetry $\to$ derived budget (Theorem~\ref{thm:budget}). Strictly weaker. \\
\addlinespace
Mild Defect Aversion (SB-1) & \textbf{Integrated} & Non-adversarial selection. Introduces no probability measure. \\
\addlinespace
Derived Locality & \textbf{Integrated} & Phase-A basins exclude global moves (Theorem~\ref{thm:locality}). One-way: Phase-A $\Rightarrow$ locality; no converse claimed. \\
\addlinespace
Separability Hierarchy & \textbf{Integrated} & Full local separability $\to$ LS-0 (probe completeness) + LS-2 (pairwise obstruction sufficiency). LS-2 verified on toy substrates. \\
\addlinespace
Linear Obstruction Growth & \textbf{Integrated} & $f(L)$ heuristic $\to$ strict $c_0(L-1)$ bound. No genericity or density argument. \\
\addlinespace
No Probe Starvation & \textbf{Integrated} & Probe richness folded into Phase-A as minimal coverage condition. \\
\addlinespace
Phase-A Failure Gates & \textbf{Integrated} & F1--F4 falsifiability conditions. Phase-A becomes modular/testable. \\
\addlinespace
Necessity Lemmas (D.7--D.10) & \textbf{Integrated} & NonRed necessary (D.7); SpecAct automatic under primitivity (D.9). Sharp ladder: $\phi_3 > 0$ and NonRed (given primitivity). \\
\addlinespace
Symmetry--NonRed Analysis (D.11) & \textbf{Integrated} & $g_2$-fiber splitting across orbits. NonRed structurally rare under symmetry. \\
\addlinespace
Certified Defect Floor & \textbf{Integrated} & $D_{\mathrm{floor}}(U) = c_\star N_{\mathrm{cert}}/T_{\mathrm{cov}}$. Structural lower bound replaces $D_{\max}$. $\tau$-selection now compares obstruction pressure vs.\ floor + slack. No dynamics in inequality. \\
\addlinespace
Recurrence Certificate & \textbf{Integrated} & Coverage epoch $T_{\mathrm{cov}}(U)$: minimal word length covering all probes. Finite combinatorial property, no ergodicity. \\
\addlinespace
Certified Phase Definitions & \textbf{Integrated} & Phase A/B/C defined by LS-2 cert.\ + recurrence cert.\ + floor small + slack small. No dynamics, temperature, or equilibrium. Pre-adaptation demoted to one sufficient mechanism. \\
\addlinespace
Admissible Words & \textbf{Integrated} & ``Histories'' reinterpreted as finite words in admissible semigroup. Eliminates temporal reading. All results algebraic. \\
\addlinespace
V7.5 Enumeration Protocol (Section~\ref{sec:v75-enum}) & \textbf{Integrated} & Canonical E/O protocol for dimension-threshold claims with explicit confidence classes and script hooks; prevents unsupported threshold assertions. \\
\addlinespace
LS-2 Certificate Format & \textbf{Open} & Machine-checkable obstruction certificates for canonical probes. \\
\bottomrule
\end{tabular}
\end{center}

\end{document}
